\documentclass[12pt,a4paper]{report}
% RoboFetch - Software Engineering Report
\usepackage[utf8]{inputenc}
\usepackage[T1]{fontenc}
\usepackage{mathpazo}
\usepackage{inconsolata}
\usepackage[a4paper,margin=2.5cm]{geometry}
\usepackage{setspace}
\usepackage{graphicx}
\graphicspath{{Figures/}{figures/}}
\usepackage{xcolor}
\usepackage{booktabs}
\usepackage{longtable}
\usepackage{tabularx}
\usepackage{array}
\usepackage{enumitem}
\usepackage{caption}
\usepackage{float}
\usepackage{fancyhdr}
\usepackage{hyperref}
\usepackage{listings}
\usepackage{titlesec}
\usepackage{amsmath}
\usepackage{amssymb}
\setlength{\parindent}{0pt}
\setlength{\parskip}{4pt}
\setlist{itemsep=2pt, topsep=4pt}
\setlist[enumerate]{leftmargin=1.6em}
\renewcommand{\arraystretch}{1.25}
\definecolor{codegreen}{rgb}{0.2,0.55,0.2}
\definecolor{codegray}{rgb}{0.45,0.45,0.45}
\definecolor{codepurple}{rgb}{0.5,0.0,0.7}
\definecolor{backcolour}{rgb}{0.97,0.97,0.98}
\definecolor{framerule}{rgb}{0.78,0.78,0.82}
\hypersetup{
    colorlinks=true,
    linkcolor=black,
    urlcolor=black,
    citecolor=black,
    pdftitle={RoboFetch Software Engineering Report},
    pdfauthor={Arian Shafagh}
}
\pagestyle{fancy}
\fancyhf{}
\fancyhead[L]{\small\itshape RoboFetch}
\fancyhead[R]{\small\itshape University of Messina}
\fancyfoot[C]{\thepage}
\setlength{\headheight}{15pt}
\renewcommand{\headrulewidth}{0.4pt}
\renewcommand{\footrulewidth}{0pt}
\titleformat{\chapter}
  {\normalfont\huge\bfseries}
  {\thechapter}{1em}{}
\titlespacing*{\chapter}{0pt}{-20pt}{22pt}
\lstdefinestyle{report}{
    backgroundcolor=\color{backcolour},
    commentstyle=\color{codegreen},
    keywordstyle=\color{codepurple},
    numberstyle=\tiny\color{codegray},
    stringstyle=\color{codegreen},
    basicstyle=\ttfamily\footnotesize,
    breakatwhitespace=false,
    breaklines=true,
    captionpos=b,
    keepspaces=true,
    numbers=left,
    numbersep=6pt,
    showspaces=false,
    showstringspaces=false,
    showtabs=false,
    tabsize=4,
    frame=single,
    rulecolor=\color{framerule}
}
\lstset{style=report}
% Shorthand for inline code so underscores never need escaping by hand.
\newcommand{\code}[1]{\texttt{#1}}
% Figure helper, so every figure is inserted the same way.
\newcommand{\uploadedfigure}[4]{
\begin{figure}[H]
    \centering
    \includegraphics[width=#2]{#1}
    \caption{#3}
    \label{#4}
\end{figure}
}
% A screenshot that has not been taken yet renders as a labelled box rather than breaking the
% build, so the document always compiles and every pending slot is visible in the PDF. Drop the
% named PNG into Figures/ and it silently becomes a real figure - nothing else to change.
\newcommand{\screenshotslot}[4]{%
\begin{figure}[H]
    \centering
    \IfFileExists{#1}{\includegraphics[width=#2]{#1}}{%
      \fbox{\parbox[c][0.26\textheight][c]{#2}{\centering\ttfamily
      SCREENSHOT PENDING\\[0.6em]\small #1\\[0.6em]\normalfont\itshape #3}}}
    \caption{#3}
    \label{#4}
\end{figure}
}
\begin{document}

% ============================================================================
% TITLE PAGE
% ============================================================================
\begin{titlepage}
\centering
\thispagestyle{empty}

% University logo. Drop unime-logo.png into Figures/ (or beside this file) and it appears
% automatically; if the file is missing the page still compiles without it.
%
% Both locations are tested on purpose. \graphicspath tells \includegraphics to look in
% Figures/, but \IfFileExists does NOT consult it - it resolves relative to the directory being
% compiled. Testing only the bare name would therefore fail for a logo sitting in Figures/, and
% the page would compile silently without ever showing it.
\newcommand{\unimelogo}{\includegraphics[width=0.28\textwidth]{unime-logo.png}\par\vspace{1cm}}
\IfFileExists{Figures/unime-logo.png}{\unimelogo}{%
  \IfFileExists{unime-logo.png}{\unimelogo}{\vspace{0.5cm}}}

{\scshape\LARGE University of Messina \par}
\vspace{0.4cm}
{\scshape\large Department of Engineering \par}   % <-- change to your department
\vspace{0.3cm}
{\large Bachelor's Degree in Computer Engineering \par}  % <-- change to your programme

\vspace{2.2cm}

{\rule{\textwidth}{0.6pt}\par}
\vspace{0.7cm}
{\huge\bfseries RoboFetch \par}
\vspace{0.5cm}
{\Large A Software Engineering Report \par}
\vspace{0.7cm}
{\rule{\textwidth}{0.6pt}\par}

\vspace{2.5cm}

% ---- Supervisor (left) and Candidate (right) -------------------------------
\begin{minipage}[t]{0.48\textwidth}
    \raggedright
    {\large\bfseries Supervisor\par}
    \vspace{0.35cm}
    {\large Prof.\ Salvatore Distefano\par}
\end{minipage}
\hfill
\begin{minipage}[t]{0.48\textwidth}
    \raggedleft
    {\large\bfseries Candidate\par}
    \vspace{0.35cm}
    {\large Arian Shafagh\par}
    {\normalsize Mat.\ XXXXXXX\par}   % <-- your matriculation number, or delete this line
\end{minipage}

\vfill

{\large Academic Year 2025/2026 \par}   % <-- adjust if needed

\end{titlepage}

% ============================================================================
% FRONT MATTER
% ============================================================================
\pagenumbering{roman}
\tableofcontents
\clearpage
\pagenumbering{arabic}

% ============================================================================

\chapter{Introduction}
% ============================================================

\section{Project Overview}
RoboFetch is a simulated warehouse robot that collects a product from a shelf and delivers it to
a delivery bay. A client opens a web page, chooses a product from a catalogue and a destination,
and the system replies with the predicted cost of the job and a decision: the robot either accepts
the order or refuses it with a reason. If the order is accepted, a differential-drive robot inside
a Gazebo simulation drives to the shelf, attaches the parcel, carries it to the bay, releases it,
and returns to its charging station.

The system is built on ROS~2 Jazzy and Gazebo Harmonic for the robotic part, and on FastAPI with
SQLite for the web part. A small machine-learning service predicts whether an order can be
completed. The whole application runs on a single machine under WSL2, and the user interface is an
ordinary web page that works in any browser.

The project is not only about making the robot move. Making a robot drive from one point to
another is largely solved by the navigation library that the project uses. What makes this a
software engineering project is everything around that: a database that describes the warehouse,
a model of the robot's physical condition, a decision procedure that decides whether a job is
possible before committing to it, a web interface, and a test strategy that can prove the robot
really did what the software claims.

\uploadedfigure{Figures/fig-warehouse-layout.png}{0.92\textwidth}
{The simulated warehouse, drawn to scale from the world description file. Three shelves hold six
products, two delivery bays sit in opposite corners, and the charging station is separate from
both.}{fig:warehouse}

\section{Project Motivation}
The first motivation is realism in the ordering process. In an early version of this project the
client had to write the name of an item and the coordinates of a shelf. That is not how a warehouse
works. A warehouse has a catalogue: every product has an identifier, a description, a weight and a
storage location. Ordering should mean choosing a product, not typing coordinates. Moving the
warehouse description into a database also means the layout can change without changing any code.

The second motivation is the question that a real fleet operator actually asks, which is not
``can the robot drive there'' but ``can the robot do this job right now''. A robot has a battery
that drains, motors that heat up, and mechanical parts that wear. A heavy product costs more energy
and produces more heat than a light one. A long route costs more than a short one. If the system
accepts every order without thinking about this, the robot will eventually run out of charge in the
middle of the warehouse. The system should be able to say no.

The third motivation is evidence. During development the system once reported that three
deliveries out of three had succeeded when in fact no parcel had moved at all. Every command had
returned success, and every status message agreed with the mistake. That experience shaped the
whole testing approach of the project: a report of success is a claim, and the position of the
parcel inside the simulator is the evidence. Wherever the software makes a physical claim, the
project checks it against the simulator rather than against its own log messages.

The fourth motivation is simplicity of the interface. The web application is rendered entirely on
the server and contains no JavaScript. This was a deliberate constraint. It keeps the client side
trivial to reason about, it removes an entire category of bugs, and it means the interface works
in any browser without a build step.

% ============================================================
\chapter{Requirements and Feature Categorization}
\label{ch:categorization}
% ============================================================

\section{Focusing on the Requirements}
The course requires a system whose functionality is balanced across three categories: features
that are oriented towards data and repositories, features that are provided by third-party
services or libraries, and complex features whose application logic is designed and implemented as
part of the project. This chapter classifies every functionality of RoboFetch into one of those
three categories and gives a short motivation for each classification.

The classification matters in this project because it contains a machine-learning model, and a
machine-learning model taken from a library is a third-party service. It is listed as such below.
The argument that the project makes is not that the model is a complex contribution, but that the
decision workflow which surrounds the model is one. That argument is made explicitly in
Section~\ref{sec:cat-complex} and demonstrated algorithmically in Chapter~\ref{ch:complex}.

\section{Data and Repository Oriented Functionalities}
These features are built on storage and simple computation. They are the record of what the
warehouse contains, what was ordered, what happened, and what condition the robot is in. There is
no novel algorithm in them; the engineering value is in the design of the schema and in keeping the
stored data honest.

\begin{longtable}{@{}p{0.06\textwidth}p{0.36\textwidth}p{0.52\textwidth}@{}}
\toprule
\textbf{ID} & \textbf{Functionality} & \textbf{Motivation for this category} \\
\midrule
\endhead
A1 & Product catalogue management & Create, read, update and delete rows describing products:
identifier, name, category, weight, stock, storage location and simulator model. \\
A2 & Warehouse layout registry & The same operations over shelves, pick points, delivery bays and
the charging station. These coordinates used to be constants inside the source code. \\
A3 & Order lifecycle management & Storing an order and moving it through its states as the robot
reports progress. Standard record management. \\
A4 & Order preview resolution & Turning a product identifier into a pick point, a route and a
distance. A lookup and a sum of distances; no decision is taken here. \\
A5 & Delivery history & An append-and-query record of duration, distance, energy, payload and
outcome for every finished order. \\
A6 & Robot telemetry time series & Periodic inserts of battery, temperature, condition, payload and
speed. \\
A7 & Per-run CSV export & Serialisation of the telemetry to a file, one file per run. \\
A8 & Statistics and key indicators & Success rate, mean delivery time, average energy and total
energy, all produced with aggregate SQL. \\
A9 & Deterministic cost formula & A closed-form expression that converts distance, payload,
temperature and condition into energy. It is a formula, not a learned or searched result, and it is
also the fallback used when the prediction service is unavailable. \\
\bottomrule
\caption{Data and repository oriented functionalities.}
\label{tab:cat-a}
\end{longtable}

\section{Third-Party Services, Libraries and Tools}
These components are integrated and configured but not implemented as part of the project.
Considerable engineering effort goes into using them correctly, but none of the logic inside them
belongs to the project.

\begin{longtable}{@{}p{0.06\textwidth}p{0.30\textwidth}p{0.58\textwidth}@{}}
\toprule
\textbf{ID} & \textbf{Service} & \textbf{Motivation for this category} \\
\midrule
\endhead
B1 & Nav2 navigation stack & A complete off-the-shelf navigation system providing localization,
global planning, local control and a behaviour tree. The project supplies parameters, a map and
goals, and receives a result. \\
B2 & Gazebo Harmonic & The physics simulator, including the lidar sensor model and the detachable
joint used as a gripper. Used as provided. \\
B3 & FastAPI and Uvicorn & The web framework and the server that runs it. \\
B4 & SQLite & The database engine. \\
B5 & scikit-learn & The implementation of the random forest classifier and the surrounding
pipeline. Considered a third-party service on its own; see
Section~\ref{sec:cat-complex}. \\
B6 & pandas and NumPy & Data handling for the training set and numerical work. \\
B7 & joblib & Serialisation of the trained model so the service can load it. \\
B8 & ROS tooling & The ROS to Gazebo bridge, the robot state publisher and the URDF processor. \\
\bottomrule
\caption{Third-party services, libraries and tools.}
\label{tab:cat-b}
\end{longtable}

\section{Complex Functionalities}
\label{sec:cat-complex}
These three features carry the application logic of the system. Each of them is described
algorithmically in Chapter~\ref{ch:complex}, with the diagrams that the course requires for complex
features.

\subsection*{C1 - Robot condition model and duty cycle}
This is a dynamic model of the robot's physical state together with the policy that decides when
the robot is allowed to work. Three quantities evolve together and influence each other. Energy
consumption grows with payload mass and distance. Temperature grows with motor load and payload and
decays towards the ambient temperature. Mechanical condition degrades with accumulated energy and
with time spent above a temperature threshold, and a degraded condition then increases energy
consumption again. That mutual influence is what makes it a model rather than three independent
counters. On top of those update rules sits a state machine which returns the robot to its station after
a period without work, refuses to depart without enough charge to come back, and forces a cooldown
when the motors overheat.

\emph{Motivation:} there is no library that models this robot. The update rules, the coupling
between them and the operational policy were all specified and implemented in this project, and the
behaviour is described by a state diagram and a set of step-by-step update rules.

\subsection*{C2 - Order admission and cost preview}
This is the procedure that converts a request into a decision. It resolves the product to its pick
point, computes the three legs of the route including the return to the station, applies the cost
formula to each leg, simulates the condition model forward to find the peak temperature the job
would produce, assembles a feature vector from live telemetry, consults the classifier, and then
applies a policy which can overrule the classifier's answer.

\emph{Motivation:} this is the classification that the course brief warns about, so it is worth
being precise. A classifier taken from a library is a third-party service, and it is listed as B5
above. What belongs to this project is the workflow around it. The model never sees an order; it
sees a vector of five numbers that this system constructs. Its answer is one input among several,
and it is consulted last, after four physical checks have already passed. A classifier that answers
``feasible'' still produces a refusal if the robot could not physically return to its station
afterwards. When the prediction service cannot be reached at all, the workflow still produces a
decision using the deterministic model alone, and records that it did so. The feature construction,
the route computation, the forward simulation, the gate ordering and the fallback are all
implemented in this project.

\subsection*{C3 - Physical delivery verification}
Before an order is marked complete, the parcel's actual position inside the simulator is compared
with the coordinates of the delivery bay. If they disagree, the order fails, whatever the commands
reported.

\emph{Motivation:} this is not a status check but a cross-validation between two independent
sources of truth, the command pipeline and the physics engine. It exists because of a real failure
during development in which the system reported three successful deliveries while no parcel had
moved. The logic and, more importantly, the decision about what counts as admissible evidence
belong to this project.

% ============================================================
\chapter{Software Development Process}
\label{ch:process}
% ============================================================

\section{Choice of Process Model}
The project follows the \textbf{incremental model}. It is a plan-driven process: the system is
delivered as a sequence of increments, and each increment runs the full sequence of phases --
requirements, analysis, design, implementation, testing and deployment -- ending in a working,
verified subset of the system rather than in a document. Pressman describes it as combining linear
and parallel process flows, each increment applying the linear sequence; Sommerville calls the same
shape incremental delivery, and McConnell staged delivery. It is one model, not two placed side by
side, and this chapter and Section~\ref{sec:increments} are organised to match it.

The single claim that distinguishes it from a waterfall is this: \textbf{requirements are frozen per
increment, not for the project as a whole.} Each increment is specified before it is built and is
traceable from its requirements to its tests, which is the discipline a plan-driven process asks
for. Between increments the specification is allowed to grow, and in this project it did: the
approved proposal defines fifteen functional requirements and this report defines twenty-one. Under
a pure waterfall those six additions would be a defect in the process. Under the incremental model
they are the process behaving as designed.

\subsection*{Why the phases are still followed inside each increment}
Three properties of this problem make the plan-driven half appropriate.

First, within an increment the requirements are knowable in advance. The requirements of a delivery
robot are determined largely by physics and by the warehouse, and both are known before the work
starts; they are not discovered by negotiation. Second, the client specifies once per increment
rather than collaborating continuously, so there is no channel through which a mid-increment change
of scope would arrive. Third, the architecture is load-bearing and expensive to revise: the
boundaries between the web tier, the robotic middleware and the simulator determine what can be
tested and how. That architecture was settled in the first increment and has not been re-cut since,
which is what allowed the later increments to add behaviour rather than rebuild structure.

\subsection*{Why delivery is nevertheless incremental}
A single-delivery waterfall would require the whole system to be built before any of it could be
demonstrated, and for this system that is not merely inconvenient but impossible to verify. The
system is physical. A claim about it can only be believed once a real robot has moved a real parcel
in the simulator, and that check cannot be deferred to a testing phase at the end, because each
layer depends on the one beneath it behaving correctly in the world rather than on paper. Chapter
\ref{ch:testing} records a case where every command reported success while no parcel had moved;
that failure is the reason verification is placed inside every increment rather than after all of
them.

Incremental delivery is also what made two later decisions possible. A delivered increment can be
judged on its behaviour rather than on its design, and two features were \emph{withdrawn} on that
evidence in the second increment: the nearest-neighbour scheduler, because serving orders by
proximity allows a customer to be overtaken indefinitely because their goods are further away, and
the retry state machine, because its complexity was not repaid by what it did. Removing a delivered
feature is a decision a single-delivery process has no natural moment for. It is the clearest
evidence in this project that the increments were real rather than a label applied afterwards.

\subsection*{Why not an agile process}
The distinction that matters is between incremental \emph{delivery} and incremental
\emph{specification}. This project delivers incrementally but specifies once per increment. An
agile process assumes the requirements are discovered through repeated contact with users, and it
answers that assumption with artifacts built for negotiation: a product backlog reordered between
iterations, stories re-estimated as understanding changes, and a velocity measure used to forecast
the next iteration.

None of those has anything to act on here, because there is no continuous channel through which the
requirements change. Adopting them would have produced ceremony without information. For the same
reason this report presents no burndown chart: none was kept, because there was no iteration
cadence to burn down against, and drawing one after the fact would be an invention rather than a
record.

\uploadedfigure{Figures/fig-incremental-model.png}{0.95\textwidth}
{The process as followed. Each increment runs the full sequence of phases and ends in a delivered,
verified subset of the system; the specification is fixed for the increment being built, not for
the project as a whole. Maintenance is drawn as a band rather than as a phase box, because it is
not bounded by an increment: it begins the moment the first increment is delivered and continues
across all of them.}{fig:incremental}

\section{Phases and Their Deliverables}
The table below gives the phases and what each produced. Every phase runs inside \textbf{every}
increment, not once for the project: each increment is specified, designed, built, tested and
delivered before the next is specified. The deliverables are listed once because their kind is the
same each time; what differs is their content, which is given per increment in
Section~\ref{sec:increments}.

\begin{longtable}{@{}p{0.22\textwidth}p{0.72\textwidth}@{}}
\toprule
\textbf{Phase} & \textbf{Deliverables} \\
\midrule
\endhead
Requirements & The user and system requirements listed in Chapter~\ref{ch:requirements}, the
actors, and the use case model. Fixed for the increment being built, and extended between
increments. \\
Analysis & The domain model, the database design, and the categorisation of functionality. \\
Design & The architecture, the component decomposition, the interfaces between components, and
the algorithms for the complex features. \\
Implementation & The nine software packages, the simulated world, and the machine-learning
pipeline, delivered across the three increments of Section~\ref{sec:increments}. \\
Testing & The automated test suites, and the acceptance suite that runs against the live system.
Each increment was verified against the simulator before the next was specified; the measured
results are reported in Chapter~\ref{ch:testing}. \\
Deployment & The launch configuration and the run instructions. Runs once per increment: each one
ends in a system that starts from a single command and is verified. \\
Maintenance & The known limitations, the problem log and the handover documentation. Unlike the
others this is \emph{not} bounded by an increment; it begins at the first delivery and continues
across every increment after it. See Section~\ref{sec:maintenance}. \\
\bottomrule
\caption{Phases of the process and what each produced.}
\label{tab:phases}
\end{longtable}

\section{The Increments}
\label{sec:increments}
Three increments were delivered. Each ran the phases described above and ended in a system that
ran and was verified before the next was specified. This section gives, for each one, what it was
asked to do, what was designed and built, and the behaviour that closed it.

\subsection{Increment 1 --- Core delivery system}
\textbf{Requirements.} FR1--FR3, FR7--FR9 and FR13: a client can order a product by name, the robot
fetches it and puts it down at a bay, a failed grab is bounded rather than endless, a delivery is
confirmed against the simulator, and the history can be read back.

\textbf{Design and construction.} The architecture that all later work rests on was settled here:
the simulated warehouse and the robot model, the map and the navigation stack, the gripper, the
task manager that chains navigation and manipulation into a delivery, and the web tier that carries
an order from a browser into the robotic system. Construction proceeded in dependency order,
because nothing above a layer can be verified until the layer beneath it works in the world:

\begin{longtable}{@{}p{0.05\textwidth}p{0.34\textwidth}p{0.55\textwidth}@{}}
\toprule
\textbf{\#} & \textbf{Stage} & \textbf{Behaviour that closed it} \\
\midrule
\endhead
1 & Simulated world and a drivable robot & Sensor, odometry, transform and velocity topics all
carried; the robot drives under manual command \\
2 & Map and autonomous navigation & Goals across the room and back both succeed; localization error
0.05 to 0.14 m \\
3 & Manipulation & Grab, release and re-attach verified; out-of-reach grabs refused \\
4 & One complete order & A full navigate, grab, carry, release sequence completes \\
5 & Order sequencing & Orders served in a different sequence from the one submitted, as designed \\
6 & Failure handling & Three attempts, then the order fails and the queue continues \\
7 & Web tier and persistence & An order submitted over the network reaches the database, then the
robot, and its status returns to the database \\
8 & Operator interface & A delivery ordered from a browser completes end to end \\
9 & Bringup, tests and documentation & The whole system starts from one command; unit, integration
and acceptance suites pass \\
\bottomrule
\caption{Construction stages within the first increment, in dependency order.}
\label{tab:increment1}
\end{longtable}

\textbf{What closed it.} A delivery ordered over the network completed, and the parcel's position
read back from the simulator agreed with the bay.

\textbf{A note on its specification.} This increment established the core product, and its detailed
requirements were settled as it was built rather than approved beforehand. That is the ordinary
role of a first increment in this model: it produces both a working core and the understanding from
which the next increment is specified. The proposal that governs the second increment is written
explicitly as a comparison against this one.

\subsection{Increment 2 --- Decision layer}
\textbf{Requirements.} FR4--FR6, FR10--FR12, FR15 and FR16, specified in the approved proposal
before any of this increment was built. The question the first increment could not answer is the
one a warehouse actually asks: not \emph{can the robot carry this}, but \emph{can it carry this
now, and what will it cost}.

\textbf{Design and construction.} Orders move from hard-coded item names to a product catalogue held
in the database, so coordinates leave the code entirely. The robot gains a model of its own
condition --- battery, temperature and mechanical wear, coupled to one another --- and every order
passes through an admission workflow that estimates its cost, consults a classifier and applies
policy gates that can overrule it. The interface is rebuilt as server-rendered pages with no
JavaScript. These are the complex functionalities, and they are described algorithmically in
Chapter~\ref{ch:complex}.

\textbf{What closed it.} Orders are refused with reasons the operator can read, and the acceptance
suite passes against the live simulator.

\textbf{Two features were withdrawn.} The nearest-neighbour scheduler and the retry state machine
were both delivered and verified in the first increment, and both were removed in this one. The
scheduler was withdrawn because watching it work made the wrong behaviour plain: serving by
proximity lets a customer be overtaken indefinitely because their goods are further away. That is a
requirements judgement which only became available once the behaviour could be observed, and it is
now fixed as FR6. The state machine was withdrawn because its structure was more elaborate than the
behaviour it produced; the guarantee it existed to provide is stated directly by FR8 and met by a
bounded loop. Both removals reduced the amount of code in the system, and the project records that
as a result rather than a loss.

\subsection{Increment 3 --- Operational control}
\textbf{Requirements.} FR14 and FR17--FR21. These did not come from rethinking the design. They came
from operating the system, which made three absences obvious: there was no way to sign in, so every
visitor had the same reach; there was no way to stop the robot once it had started; and there was no
way to ask it to go home. The same period added the reservation of work already accepted, without
which an admission decision was being taken against a robot that was standing still only on paper.

\textbf{Design and construction.} Authentication with two roles enforced on every request; a
non-latching emergency stop handled on the subscription callback, because the worker thread is
normally blocked waiting on a navigation result when the button is pressed; a return to the station
queued as an ordinary task so that it cannot interrupt a carry; and a reservation ledger derived
from the orders table rather than stored beside it, so that a terminal status releases the
reservation automatically and the two cannot drift apart.

\textbf{What closed it.} The acceptance suite covers all four: the login gate on every protected
page, the stop verified against the robot's measured displacement in the simulator, the queued
return, and the reservation both held and released.

\textbf{Why this is an increment and not a redesign.} None of it changed the shape of the system.
Each is a page, a queue entry or a database query fitted to machinery that already existed, which
is why the work was added to the existing architecture rather than prompting a new one.


\section{Maintenance}
\label{sec:maintenance}
Maintenance is the phase that does not fit inside an increment, and saying where it belongs is part
of describing this process honestly.

In a single-delivery process the answer is simple: maintenance is what happens after the one
delivery. Here there are three deliveries, and the moment the first increment was delivered, every
subsequent change to it became maintenance by definition. The second increment did not only add
code to the first: it \emph{modified delivered code}, removing the scheduler, replacing the retry
state machine, changing the order of service to arrival order and rebuilding the interface. That
work is maintenance carried out as part of the next increment. Sommerville makes the general point
that in iterative processes the distinction between development and maintenance stops being a
useful one, and this project is an instance of it.

Maintenance is therefore drawn in Figure~\ref{fig:incremental} as a band beginning at the first
delivery and running past the last increment, rather than as a box at the end of a row.

\subsection*{The three kinds, and where each occurred}
Maintenance is classified here as corrective, adaptive and perfective, following Swanson's
categories. All three occurred, and each is traceable to specific work.

\begin{longtable}{@{}p{0.16\textwidth}p{0.78\textwidth}@{}}
\toprule
\textbf{Kind} & \textbf{Where it occurred in this project} \\
\midrule
\endhead
Corrective &
Repairing faults in delivered behaviour. A stale web process kept answering on the same port while
the real one failed to start, so orders were accepted and never executed. The stop script did not
match the condition monitor's process name, so five of them accumulated and the web tier read
telemetry from several stale robots at once. An acceptance check counted a refused order as one
that had started, and so reported a scheduling fault that did not exist. \\
Adaptive &
Adjusting to a changed environment rather than to a changed requirement. The web framework
deprecated the startup-event interface the application used, which was migrated to its replacement.
The simulator's predecessor reached end of life and its attachment mechanism no longer existed on
the supported platform, so the gripper was rebuilt on the detachable joint the current simulator
provides. \\
Perfective &
Improving what already worked, in behaviour or in maintainability, without any fault having been
reported. The battery capacity was retuned twice so that admission control had something real to
refuse, each change requiring the classifier to be retrained because the training label depends on
it. The controller's sampling parameters were reduced when they starved the control loop, and the
pose publication rate was lowered when it saturated a node's executor. The templates were
reorganised and the test suite reduced to a representative set. Several changes were made to reduce
the chance of a fault rather than to remove one: the stop script gained a self-audit that names any
running process it would fail to kill, so the next omission is reported rather than discovered; a
test fails the build if guidance text reappears on the web pages, enforcing a rule that would
otherwise erode; and the acceptance table in this report is generated from the recorded results
rather than transcribed, after a transcribed copy was found to disagree with the figure printed
beside it. \\
\bottomrule
\caption{The three kinds of maintenance and where each occurred.}
\label{tab:maintenance}
\end{longtable}

\subsection*{What makes the maintenance possible}
Two things carry most of the weight, and both were built during the increments rather than added
afterwards.

The first is the regression suite. Because every increment ended in a verified system, a change
made later can be checked against the same criteria that closed the increment it touches. Nothing
in the maintenance record above was accepted on the strength of an argument that it ought to work.

The second is the record of what was tried and rejected. Several of the corrections above were
first attempted in ways that failed, and those attempts are documented alongside the fixes, so that
a later reader does not repeat them. Four attempts to improve delivery accuracy were reverted before
the cause was correctly identified at the grab rather than at the release; the dead helper from one
of them is still in the source, deliberately unused, with a comment explaining why calling it makes
the problem worse. A maintenance record that lists only what worked is of limited use to whoever
maintains the system next.

\section{Traceability}
Every requirement in Chapter~\ref{ch:requirements} maps forward to a design element and to at
least one test. The mapping is given in Appendix~\ref{app:traceability}. Traceability is what makes
a plan-driven process verifiable: because the requirements were fixed before design started, it is
possible to check at the end that each one is satisfied by something concrete and demonstrated by
something measurable.

% ============================================================
\chapter{Requirements Engineering}
\label{ch:requirements}
% ============================================================

\section{Levels of Requirement}
Requirements are recorded at two levels, because they answer two different questions and are read
by two different audiences.

\textbf{User requirements} state what the warehouse needs, in the language of the people who run
it. They are deliberately free of any mention of how the software works: no tables, no endpoints,
no classifiers. Each is written as a user story so that the actor, the need and the reason are all
visible in one line.

\textbf{System requirements} state what the software must do in order to satisfy a user
requirement. They are the specification the implementation is checked against, and each one names
the user requirement it refines. They are given in Section~\ref{sec:sysreq}, divided into
functional requirements, which describe behaviour, and non-functional requirements, which constrain
how that behaviour is delivered.

Every requirement carries a unique identifier that is used unchanged throughout this report, in the
traceability matrix in Appendix~\ref{app:traceability}, and in the acceptance suite itself, where
the identifier is printed beside each result. A requirement, its design and its evidence can
therefore always be lined up.

\section{User Requirements}

\begin{longtable}{@{}p{0.07\textwidth}p{0.72\textwidth}p{0.15\textwidth}@{}}
\toprule
\textbf{ID} & \textbf{User requirement} & \textbf{Refined by} \\
\midrule
\endhead
UR1 & As a warehouse controller, I want to have a product brought to a delivery bay by naming the
product, so that I do not need to know where anything is stored. & FR1, FR2, FR3, FR7 \\
UR2 & As a controller, I want to see what an order will cost and whether it is possible before I
commit to it, so that I can decide with the facts in front of me. & FR4 \\
UR3 & As a warehouse manager, I want the system to refuse work it cannot finish and tell me why,
so that the robot is not sent out on jobs that will strand it. & FR5, FR17 \\
UR4 & As a customer, I want orders served in the sequence they were placed, so that nobody is
overtaken because their goods are further away. & FR6 \\
UR5 & As a warehouse manager, I want a delivery reported as complete only when the goods are
physically at the bay, so that the records describe the warehouse and not the software. & FR8, FR9 \\
UR6 & As a controller, I want to see the robot's condition and trust that it will not be worked
past its limits, so that it is not damaged by the work I give it. & FR10, FR11 \\
UR7 & As a warehouse manager, I want to review what has happened and how well the system performs,
so that I can account for the operation. & FR13, FR15 \\
UR8 & As a controller, I want stock to stay truthful, so that I cannot order something that has
already been taken off its shelf. & FR16 \\
UR9 & As a controller, I want to stop the robot at once in an emergency and to send it home when I
choose, so that I keep direct control of a machine that is moving. & FR18, FR19 \\
UR10 & As a warehouse manager, I want only authorised people to operate the system, with
administration separated from operation, so that routine use cannot change the catalogue. & FR14,
FR20, FR21 \\
UR11 & As a warehouse manager, I want the robot's behaviour recorded, so that performance can be
analysed and improved over time. & FR12 \\
\bottomrule
\caption{User requirements, expressed as user stories, and the system requirements that refine
them.}
\label{tab:ur}
\end{longtable}

\section{System Requirements}
\label{sec:sysreq}
The system requirements below refine the user requirements of the previous section into statements
about the software itself. Where a user requirement says what the warehouse needs, a system
requirement says what the software must do to provide it, in terms precise enough to be built
against and checked.

They are given in two groups. The \emph{functional} requirements describe what the system does;
the \emph{non-functional} requirements constrain how it does it. Every one of them names its
parent user requirement, so the chain from a need in the warehouse to a line in the acceptance
suite is unbroken in both directions.

\subsection{Functional System Requirements}
Each functional requirement names the user requirement it serves and the criterion by which it is
accepted. The criterion is what makes the requirement testable: it states the observable outcome
that decides whether the requirement is met, and it is the same condition the corresponding check
asserts.

\begin{longtable}{@{}p{0.06\textwidth}p{0.05\textwidth}p{0.44\textwidth}p{0.37\textwidth}@{}}
\toprule
\textbf{ID} & \textbf{UR} & \textbf{Requirement} & \textbf{Acceptance criterion} \\
\midrule
\endhead
FR1 & UR1 & The system shall maintain a catalogue of products, each with a unique identifier, name,
category, weight and storage location. & The catalogue endpoint returns all six products with
weights and a resolved pick point for each. \\
FR2 & UR1 & The system shall maintain a description of the warehouse layout: shelves, per-product
pick points, delivery bays and one charging station. & The layout endpoint returns two delivery
bays, and every product resolves to non-null pick coordinates. \\
FR3 & UR1 & A client shall submit an order by product identifier and destination, without knowing
any coordinates. & An order posted with only a product identifier and a bay name is accepted, and
the coordinates are resolved by the system. \\
FR4 & UR2 & The system shall present a preview before the order is confirmed, containing the
product details, the route distance, the estimated energy and time, and the decision. & A preview
returns route distance, energy, peak temperature and a verdict, and creates no order. \\
FR5 & UR3 & The system shall refuse an order it predicts the robot cannot complete, and state a
reason a human can read. & An order for an item that is out of stock is refused, and the response
carries a reason naming the cause. \\
FR6 & UR4 & Accepted orders shall be executed in the order in which they were placed. & Two orders
queued in sequence begin execution in that same sequence. \\
FR7 & UR1 & The robot shall navigate to the pick point, attach the parcel, carry it to the delivery
bay and release it. & An accepted order reaches the completed state having passed through
navigating, grabbing, delivering and releasing. \\
FR8 & UR5 & A failed attachment shall be retried up to three times; after the third failure the
order fails and the robot returns to its station. & An order whose parcel cannot be reached fails
after exactly three attempts, with the attempt count recorded, and the robot returns to the
station. \\
FR9 & UR5 & The system shall confirm a delivery against the parcel's real position in the simulator
before marking the order complete. & The parcel's pose read from the simulator lies within the
delivery tolerance of 0.7 m of the bay; otherwise the order is failed. \\
FR10 & UR6 & The system shall continuously track the robot's battery, temperature, condition,
payload and cumulative distance. & All five quantities are readable from the robot endpoint and
change as work is carried out. \\
FR11 & UR6 & After a period without pending orders the robot shall return to its station and
recharge. & With no work queued, the robot reaches the station, reports itself charging, and its
battery rises. \\
FR12 & UR11 & Every run shall be logged to a file suitable for machine-learning training. & A
per-run file exists after a delivery and contains at least one sample row per recorded interval. \\
FR13 & UR7 & A client shall be able to view the history of all orders and their outcomes. & The
history lists every order of the session with its terminal state and detail. \\
FR14 & UR10 & An administrator shall be able to create, read, update and delete products and layout
entries. & A product can be created, renamed, re-weighted and deleted through the application, and
each change is visible afterwards. \\
FR15 & UR7 & The system shall report aggregate statistics: success rate, mean delivery time, energy
per delivery and totals. & The reported totals equal the totals counted directly from the orders
table. \\
FR16 & UR8 & A product that has been delivered shall be removed from stock and shall stop being
orderable. & After a delivery the product is absent from the orderable list, remains in the full
catalogue at zero stock, and a further order for it is refused. \\
FR17 & UR3 & An accepted order shall reserve its projected energy and peak temperature until it
reaches a terminal state, and the next decision shall be taken against what would remain. & While
an order is in flight the next preview is costed from the reduced battery; when the order ends the
reservation is released and no energy remains held. \\
FR18 & UR9 & A controller shall be able to queue a return to the station. It shall be costed and
reserved like any other task, and shall never be refused. & A queued return waits for the work
ahead of it, then runs, and the robot ends docked and charging. \\
FR19 & UR9 & A controller shall be able to halt the robot immediately, failing the running order
and everything queued. The halt shall not persist once it has been carried out. & After the stop
the robot's measured displacement in the simulator is zero, the running and queued orders are
failed, and the very next order is accepted normally. \\
FR20 & UR10 & Every page shall require an authenticated session, and the two roles shall differ in
what they may reach, enforced on the request. & Each of the eight protected pages redirects an
anonymous visitor to the login page; a controller is refused the administration pages that an
administrator reaches. \\
FR21 & UR10 & An administrator shall be able to read any database table through the application,
with password hashes, salts and session tokens redacted. & A table name that is not in the schema
is rejected, and every redacted column reads as redacted rather than as its stored value. \\
\bottomrule
\caption{Functional requirements, the user requirement each serves, and the criterion by which each
is accepted.}
\label{tab:fr}
\end{longtable}

\subsection{Non-Functional System Requirements}
A non-functional requirement that cannot be measured cannot be enforced, so each one below states
the quantity, the method of measurement and the threshold that decides it. Writing these criteria
was not a formality: doing so revealed that NFR4 was not being met, which is recorded in
Section~\ref{sec:nfr4} and in the limitations.

\begin{longtable}{@{}p{0.06\textwidth}p{0.40\textwidth}p{0.46\textwidth}@{}}
\toprule
\textbf{ID} & \textbf{Requirement} & \textbf{How it is measured, and the threshold} \\
\midrule
\endhead
NFR1 & The order preview, including the prediction, shall be returned within two seconds. & Wall
clock time around a preview request that reaches the prediction service, measured by the acceptance
suite. Threshold: under 2000 ms. Measured: 416 ms. \\
NFR2 & The system shall remain operable when the prediction service is unavailable, falling back to
a deterministic rule. & With the prediction client replaced by one that always fails, a preview
still returns a decision, and the response records that the decision was taken without the model.
Threshold: a decision is produced and is attributed to the policy. \\
NFR3 & The web application shall require no JavaScript. & Every protected page is fetched and its
body searched for a script element. Threshold: all eight pages return successfully and none
contains one. \\
NFR4 & Robot telemetry shall be recorded at least once per second of simulated time while the robot
is active. & The largest gap between consecutive telemetry rows over a delivery, taken from the
recorded timestamps. Threshold: no gap greater than one second on the stated clock. See
Section~\ref{sec:nfr4}. \\
NFR5 & Every reported delivery shall be verifiable against the state of the simulator. & For each
completed order, the parcel's pose is read from the simulator and compared with the bay. Threshold:
within 0.7 m, and any order outside it is failed rather than reported complete. \\
NFR6 & The machine-learning training pipeline shall be removable without affecting the running
system. & Static check that no module under the source tree imports the training package.
Threshold: zero such imports, so the generator and trainer can be deleted. \\
\bottomrule
\caption{Non-functional requirements, each with the measurement that decides it.}
\label{tab:nfr}
\end{longtable}

\section{Actors and Use Cases}
The system has two human actors, and they correspond exactly to the two roles the login enforces.
The \emph{controller} places and tracks orders, inspects the condition of the robot, sends it back to
its station and stops it in an emergency. The \emph{administrator} may do all of that and in addition
maintains the catalogue and the layout, manages accounts, inspects the statistics and browses the
database. Requirements written below in terms of a \emph{client} refer to the controller role. The navigation stack and the simulator appear as a supporting system actor,
because the project treats them as external services rather than as parts it implements.

\uploadedfigure{Figures/01-use-case.png}{0.95\textwidth}
{Use case diagram. The preview includes the feasibility prediction, and confirming an order
includes obtaining a preview first, which is why an order can never be placed without a decision
having been made. The administrator generalises the controller rather than duplicating its arrows,
because the roles nest: an administrator may do everything a controller may.}{fig:usecase}

\begin{longtable}{@{}p{0.08\textwidth}p{0.28\textwidth}p{0.56\textwidth}@{}}
\toprule
\textbf{ID} & \textbf{Use case} & \textbf{Description} \\
\midrule
\endhead
UC1 & Browse product catalogue & The controller sees the available products with their weights and
storage locations. \\
UC2 & Request order preview & The controller selects a product and a destination and receives the
predicted cost and the decision, without committing to anything. \\
UC3 & Confirm an order & The controller accepts the preview and the order is stored and sent to the
robot. \\
UC4 & View order history & The controller sees all orders, their states and their outcomes. \\
UC5 & Execute a delivery & The robot performs the pick-and-place sequence and verifies it. \\
UC6 & View analytics & The aggregate statistics are reported: success rate, times, energy. \\
UC7 & View robot condition & Battery, temperature, condition and duty state are inspected. \\
UC8 & Browse the database & The administrator reads any table, a page at a time, with secrets
redacted. \\
UC9 & Predict order feasibility & The system decides whether an order is possible. Always invoked
as part of a preview, never on its own. \\
UC10 & Return the robot to its station & The controller queues a trip home. It waits its turn like
any other task and is never refused. \\
UC11 & Stop the robot immediately & The controller abandons all work and halts the robot in one
press, with no confirmation and no latch. \\
UC12 & Sign in & A visitor authenticates and is granted the reach of their role; an anonymous
visitor is redirected to the login page. \\
UC13 & Manage products and layout & The administrator creates, updates and deletes products and
delivery points. \\
\bottomrule
\caption{Use cases.}
\label{tab:usecases}
\end{longtable}

% ============================================================
\chapter{Analysis and Domain Model}
\label{ch:analysis}
% ============================================================

\section{Domain Concepts}
The analysis phase identified the entities that the system reasons about. A \emph{product} is
something a client can order; it has a weight, which matters because weight drives energy
consumption and heat. A \emph{shelf} is a physical structure holding products. A \emph{pick point}
is the position the robot parks at in order to reach one particular product; each product has its
own pick point, because two products on the same shelf are a metre apart and the robot cannot
reach both from the same place. A \emph{delivery bay} is a destination. The \emph{station} is where
the robot charges and waits. An \emph{order} links a product to a delivery bay and carries the
decision that was made about it. \emph{Telemetry} is the record of the robot's condition over time.

\uploadedfigure{Figures/02-domain-class.png}{0.98\textwidth}
{Domain class diagram of the runtime objects. The condition model on the left is shared: the same
module is used by the robot node, by the admission workflow and by the tool that generates
training data.}{fig:domainclass}

\section{Design of the Database}
The database is the part of the system that changed most between versions, and the design is worth
explaining because it determines how flexible the rest of the system is.

The central decision is that an \emph{order references a product, never a coordinate}. In the
first version an order carried the coordinates of a pickup point and a drop-off point. That worked,
but it meant the client had to know where things were, and it meant that moving a shelf would
invalidate every stored order. Now the client sends a product identifier and a destination
identifier, and the system resolves those to coordinates at the moment the order is admitted. The
consequence is that the warehouse can be rearranged by editing rows in a table.

Figure~\ref{fig:er} shows the schema as an entity-relationship diagram. Each box is a table, each
row inside it is a column with its type, and the connectors carry the cardinality of the
relationship, so the diagram shows both what is stored and how the tables relate to one another.

\uploadedfigure{Figures/06-er-diagram.png}{0.95\textwidth}
{The database schema. Primary keys are marked, foreign keys link the tables, and the cardinality
of each relationship is shown on the connectors.}{fig:er}

The same schema is given again in Figure~\ref{fig:dbclass} as a UML class diagram, with each table
as a class carrying the \texttt{<<table>>} stereotype and each column its SQL type. The two views
answer different questions: the entity-relationship diagram shows how the tables relate, while the
class diagram shows precisely what each one stores, which is the form needed when reading the
schema against the code that creates it.

\uploadedfigure{Figures/07-database-class.png}{0.95\textwidth}
{The database schema as a UML class diagram. Each table is a class with the table stereotype;
columns carry their SQL types, and primary and foreign keys are marked.}{fig:dbclass}

The schema contains nine tables in three groups. The layout group describes the warehouse:
\code{shelves}, \code{pick\_points}, \code{delivery\_points} and \code{station}. The catalogue and
order group contains \code{products}, \code{orders} and \code{delivery\_history}. The condition
group contains \code{robot\_telemetry} and \code{robot\_state}.

Four design decisions are worth justifying:

\textbf{The layout is data, not constants.} Every coordinate lives in a table. Nothing in the
robot code or the web code contains a hard-coded position. Adding a product or moving a delivery
bay requires no recompilation.

\textbf{The admission estimate is stored on the order.} When an order is admitted, the predicted
distance, energy and duration are written into the order row along with the decision and the reason
for it. This costs a few columns and makes it possible to compare what was predicted against what
actually happened, which is the basis of Figure~\ref{fig:predvsactual}.

\textbf{Current state and history are separate tables.} \code{robot\_state} holds a single row with
the current condition so a web page can read it with one query, while \code{robot\_telemetry} is
append-only and holds the full time series used for analysis and for machine-learning training.
Mixing the two would have made either the page slow or the history unavailable.

\textbf{Two tables intentionally have no relationships.} \code{station} and \code{robot\_state} are
single-row tables. The station is a property of the warehouse rather than of any order, and the
robot state is a property of the robot. Leaving them unconnected is correct, and the diagram
annotates this so it does not look like an omission.

The schema is in third normal form. Each fact is stored once: a product's weight lives only in the
\code{products} table, and a pick point's coordinates live only in \code{pick\_points}. Orders refer
to both by key.

% ============================================================
\chapter{Software Architecture}
\label{ch:architecture}
% ============================================================

\section{Architectural Drivers}
Architecture is the set of decisions that are expensive to change later, so it is worth stating
what actually drove them in this project. Four concerns shaped the structure more than anything
else.

\textbf{The robot must keep working when the prediction service does not.} A machine-learning
service is the least reliable component in the system: it may be stopped, it may be slow, or it may
have no model loaded. If the ability to accept an order depended on it, the whole system would
inherit its reliability. This drove the decision to place the model behind a network boundary and
to require a deterministic fallback.

\textbf{Physical claims must be independently verifiable.} The software must never be the only
witness to its own success. This drove the decision to keep a channel from the simulator that
reports the real position of every parcel, separate from the command path.

\textbf{The warehouse description must be changeable without touching code.} This drove the
database-centred design described in the previous chapter.

\textbf{The client must be simple.} Requiring no JavaScript drove the choice of server-rendered
pages and removed the need for any client-side state.

\section{Layered Architecture}
The system is organised into five layers. Each layer depends only on the layers below it; nothing
in a lower layer knows that the layers above exist. This is what allows the robot logic to be
tested without a web server and the decision logic to be tested without a robot.

\uploadedfigure{Figures/fig-layered-architecture.png}{0.85\textwidth}
{The five layers and the direction of dependencies.}{fig:layers}

The \textbf{presentation layer} contains only templates and a stylesheet. It holds no logic and no
state.

The \textbf{application layer} is the web tier. It owns the database, the admission workflow, the
client for the prediction service and the bridge to the robotic middleware. It is the only layer
that talks to more than one of its neighbours, and that is its purpose: it is the place where a
request from a person becomes a command for a machine.

The \textbf{prediction service} is a separate process holding the trained model.

The \textbf{robot logic layer} contains the software written for this robot: the task manager that
executes an order, the condition model, and the gripper node.

The \textbf{robot platform layer} contains the third-party components: the navigation stack, the
simulator, and the bridge between them.

\section{Component Decomposition}
Figure~\ref{fig:component} shows the components, their interfaces and the direction of every
dependency. Table~\ref{tab:components} states the single responsibility of each package.

\uploadedfigure{Figures/08-component.png}{0.98\textwidth}
{Component diagram showing the nine packages and their interfaces.}{fig:component}

\begin{longtable}{@{}p{0.26\textwidth}p{0.68\textwidth}@{}}
\toprule
\textbf{Package} & \textbf{Single responsibility} \\
\midrule
\endhead
\code{robofetch\_description} & Describes the robot: geometry, wheels, sensors, gripper link and
the simulator plugins attached to them. \\
\code{robofetch\_gazebo} & Describes the world and the translation between simulator topics and
middleware topics. \\
\code{robofetch\_nav} & Holds the navigation configuration and the generated map. \\
\code{robofetch\_interfaces} & Defines the one custom service used to command the gripper. \\
\code{robofetch\_core} & The robot's own logic: condition model, task manager, gripper node. \\
\code{robofetch\_bridge} & The web tier: HTTP interface, database access, admission workflow and
the client for the prediction service. \\
\code{robofetch\_ai} & Serves the trained classifier over HTTP. \\
\code{robofetch\_web} & Templates and stylesheet. \\
\code{robofetch\_bringup} & Starts every process in the correct order. \\
\bottomrule
\caption{Package responsibilities.}
\label{tab:components}
\end{longtable}

\section{Integration and Communication Styles}
The system uses three different communication styles, each chosen for a different reason.

\textbf{Publish and subscribe, between the web tier and the robot.} Orders travel from the web
process to the task manager on a topic, and status changes travel back on another topic. The
payload is JSON carried inside a plain string message. Using a generic string rather than a custom
message type keeps the web tier free of any dependency on the robotic middleware's message
definitions, which matters because the web tier is also the part that talks to the browser and the
database. The cost of this choice is that the payload is not type-checked by the middleware; the
benefit is a much cleaner boundary. The asynchronous style is also correct on its own terms: a
status update is a stream of events, and a dropped message must not block the HTTP request that
created the order.

\textbf{Request and response, between the browser and the web tier.} Ordinary HTTP, with pages
rendered on the server. A form post produces a new page.

\textbf{Request and response over a network boundary, between the web tier and the prediction
service.} The model could have been imported directly into the web process. Putting it behind HTTP
in a separate process costs a little latency and gains something more valuable: the failure of the
model becomes a visible, ordinary, testable event rather than a hypothetical one. Because the
boundary exists, the fallback path can be exercised simply by stopping a process, and the system
can report on a status page whether the service is currently reachable.

\section{Data Architecture}
Database access is confined to a single class which exposes methods in the vocabulary of the
domain, such as ``list products'' or ``create order'', rather than exposing SQL to its callers.
This is the repository pattern, and it means the rest of the system does not know that SQLite is
being used.

Each operation opens its own short-lived connection. This is unusual enough to justify: the web
process serves HTTP requests on one set of threads and receives robot messages on another, and both
write to the database. A single shared connection would need locking around every use. Short-lived
connections avoid the problem entirely, and the cost is negligible for a local file database.

\section{Deployment Architecture}
Every component runs on one machine. Figure~\ref{fig:deployment} shows the process layout.

\uploadedfigure{Figures/09-deployment.png}{0.98\textwidth}
{Deployment diagram. The browser is the only component outside the Linux environment.}
{fig:deployment}

The browser is deliberately the only component outside the Linux environment. Serving the interface
over HTTP means the user interacts with the system through an ordinary browser on the host
operating system, which avoids the graphical problems that the simulator's own windows suffer from
in this environment. It also means the interface would work unchanged from another machine on the
same network.

\section{Architectural Decisions}
Table~\ref{tab:adr} records the significant decisions, the alternatives considered and the
consequences accepted.

\begin{longtable}{@{}p{0.20\textwidth}p{0.22\textwidth}p{0.24\textwidth}p{0.24\textwidth}@{}}
\toprule
\textbf{Decision} & \textbf{Alternative} & \textbf{Rationale} & \textbf{Consequence} \\
\midrule
\endhead
Prediction model in a separate process behind HTTP & Import the model into the web process &
Makes the failure of the model an ordinary, testable event and keeps the web tier free of the
machine-learning dependency & Extra process to start; a small latency cost per preview \\
JSON on generic string messages between web tier and robot & Custom message types &
Keeps the web tier independent of middleware message definitions & Payload is not type-checked by
the middleware, so both sides must agree by convention \\
Map generated analytically from the world geometry & Build the map by driving the robot around &
Zero drift, perfectly repeatable, and the map frame equals the world frame so waypoints are world
coordinates & The map must be regenerated whenever the world geometry changes \\
Server-rendered pages with no JavaScript & A client-side application &
Simplicity, no build step, works in any browser & Live values need a page refresh rather than a
push \\
SQLite & A client-server database &
Single file, no service to administer, adequate for one robot & Would need replacing for
concurrent writers at scale \\
Layout stored in the database & Constants in source code &
The warehouse can be rearranged without touching code & An empty or wrong database produces a
system that starts but cannot route anything \\
\bottomrule
\caption{Architectural decisions and their consequences.}
\label{tab:adr}
\end{longtable}

% ============================================================
\chapter{Software Engineering Principles and Quality Standards}
\label{ch:quality}
% ============================================================

This chapter argues about the \emph{quality} of the implementation rather than about what it does.
A system can satisfy every functional requirement and still be badly built. The principles below
are the ones the project applied consciously, each with the concrete place in the system where the
effect can be seen.

\section{Separation of Concerns and Modularity}
The system is divided into nine packages, and each has one reason to change. If the warehouse
layout changes, only the world description and the database seed change. If the appearance of the
interface changes, only the templates and the stylesheet change. If the physics of the condition
model changes, only one module changes. This is the practical test of modularity: a change request
should map to a small, predictable set of files.

\section{Single Source of Truth}
The condition model exists in exactly one module. That module is imported by the robot node that
runs on the live robot, by the admission workflow in the web tier that predicts the cost of an
order, and by the tool that generates training data for the classifier.

This matters more than it appears. If the web tier had its own copy of the energy formula, its
predictions would slowly stop matching the robot's actual behaviour as one copy was edited and the
other was not. If the data generator had a third copy, the classifier would be trained on physics
that the real system does not implement. Sharing one definition makes that class of inconsistency
impossible rather than merely unlikely.

\section{High Cohesion and Loose Coupling}
Each module groups things that change together and knows as little as possible about the others.
The robot side never sees the product catalogue: it receives a model name, a weight and two pairs of
coordinates, and that is all it needs. The web tier never deals in coordinates until the moment it
resolves an order. The gripper node knows how to attach and detach a parcel but nothing about
orders or routes.

\section{Dependency Direction and Testability}
Testability is not a property that can be added at the end; it follows from the direction of
dependencies. In this system the two modules that contain the important decisions, the condition
model and the admission policy, contain no imports of the robotic middleware, no HTTP client and no
database handle. Everything they need is passed in as plain data.

\uploadedfigure{Figures/fig-dependency-direction.png}{0.88\textwidth}
{Dependency direction. The modules holding the decision logic depend on nothing, which is why they
can be tested in milliseconds.}{fig:deps}

The consequence is measurable. The complete decision path of the system, including every policy
gate and the full condition model, is covered by tests that run in a few seconds without starting a
simulator, a database or a web server. Had the policy been written directly inside the HTTP
handlers, testing a refusal would have required a running robot with a flat battery.

\section{Fail-Safe Design and Graceful Degradation}
The prediction service is optional by construction. If it cannot be reached, the client that talks
to it returns nothing rather than raising, and the admission workflow continues using the
deterministic model alone. The decision is still made, the order is still handled, and the system
records that the decision was taken without the model so the difference is visible afterwards.

This is the difference between a system that degrades and one that fails. The requirement NFR2
exists precisely to force this behaviour, and it is verified by a test that supplies an
unreachable prediction service and asserts that a decision is still produced.

\section{Defensive Verification}
The strongest principle the project adopted is that \emph{a command returning success is not
evidence that anything happened}. It was adopted after a failure in which every layer of the system
agreed that three deliveries had succeeded while no parcel had moved.

The principle is applied in two places. A delivery is confirmed by reading the parcel's real
position from the simulator and comparing it with the destination, not by observing that the
release command returned success. A grab is confirmed by reading the joint's own reported state,
which is either attached or detached, rather than by measuring the distance between the parcel and
the robot. The distance check had a blind spot that is worth describing because it is subtle: it
measured the gap before the attach and again afterwards, and accepted the grab if the gap was still
small. If the attach silently did nothing, neither the robot nor the parcel moved, so the second
measurement equalled the first and the check passed. A test that cannot fail is worse than no test,
because it produces confidence without evidence.

\section{Configuration over Hard-Coding}
Warehouse geometry, the product catalogue and the robot's operating parameters are all data. The
layout lives in the database; the robot's thresholds are declared as parameters that can be
overridden when the system starts. The rule applied throughout is that a value which a user might
reasonably want to change should not require an edit to a source file.

\section{Observability}
A system that cannot be inspected cannot be diagnosed. Three mechanisms provide this. A health
endpoint reports not only that the web service is alive but whether a robot is actually subscribed
to the order topic, which distinguishes ``the system is broken'' from ``nothing is listening''.
Telemetry is published continuously and stored. Every run writes a log file.

The health endpoint deserves a note, because it was added in response to a specific class of
confusion. Publishing a message to a topic that nobody subscribes to succeeds silently. An order
sent to a robot that is not running therefore looks exactly like an order that was accepted
normally. Reporting the live subscriber count turns an invisible failure into a visible one.

\section{Error Handling and Defensive Programming}
Input is validated at the boundary: an unknown product or destination produces a clear error
before anything is stored. HTTP status codes are used for their intended meanings, so a request to
cancel an order that has already started returns a conflict rather than a generic failure. Retries
are bounded: three attempts, then the order fails and the robot goes home, so no failure can loop
forever. Where a component can be absent, its absence is handled explicitly rather than allowed to
raise.

\section{Coding Standards and Documentation}
The code follows a consistent naming style and a consistent structure: modules begin with a
docstring that explains what the module is for and, where a decision is not obvious, why it was
made that way. Comments in this project are used mainly to record reasoning rather than to restate
the code, because the reasoning is the part that cannot be recovered by reading the implementation.
A separate handover document records the problems encountered during development and their
solutions.

\section{Maintainability and Extensibility}
The practical test of maintainability is how much has to change in order to make a plausible new
request work. Adding a product to the catalogue requires a row in a table and a model in the world
description; no code changes. Moving a delivery bay requires editing two numbers in a table.
Changing how much energy the robot uses requires editing one module, and every part of the system
that depends on it, including the training data generator, follows automatically.

\section{Summary of Quality Attributes}

\begin{longtable}{@{}p{0.24\textwidth}p{0.70\textwidth}@{}}
\toprule
\textbf{Attribute} & \textbf{How it is achieved and where it can be seen} \\
\midrule
\endhead
Reliability & Bounded retries, a fallback for the prediction service, and physical verification
before any success is reported. \\
Testability & Decision logic free of infrastructure dependencies; 16 automated tests that need no
simulator. \\
Maintainability & One reason to change per package; configuration held as data. \\
Observability & Health reporting, continuous telemetry, per-run logs. \\
Portability & The interface is a web page; the system runs from a single command. \\
Performance & The preview, including the prediction, responds well inside its two-second budget. \\
\bottomrule
\caption{Quality attributes and the mechanisms that provide them.}
\label{tab:quality}
\end{longtable}

% ============================================================
\chapter{The Simulated Environment}
\label{ch:environment}
% ============================================================

\section{The Warehouse}
The warehouse is a rectangular room of eight metres by six, containing three shelves, two delivery
bays and a charging station, shown to scale in Figure~\ref{fig:warehouse}. Almost every dimension in
it is the result of a specific problem, and this section explains the reasoning because the geometry
is part of the engineering rather than decoration.

\textbf{The room is rectangular, not square.} An early version used a square room. A square room is
rotationally ambiguous to a lidar: rotated by ninety degrees, the predicted scan is identical, so
the localization system genuinely cannot tell the two positions apart. The consequence was that the
robot's estimate of its own position would occasionally jump to a wrong hypothesis, after which it
would drive confidently into a wall. Making the room rectangular removes the ambiguity.

\textbf{The shelves have different sizes and are placed asymmetrically.} This gives every part of
the room a distinctive appearance to the laser, which keeps localization stable. The shelves are
therefore doing two jobs at once: they are the storage locations, and they are the landmarks.

\textbf{Every shelf is flush against a wall.} A gap between a shelf and a wall forms a corridor. A
robot of forty-four centimetres, with the safety margin the planner adds around obstacles, cannot
turn around in a corridor narrower than roughly 1.2 metres. During development the robot drove into
such a corridor, wedged itself, and its position estimate then drifted by nearly three metres,
after which the planner refused every goal it was given, including one in the empty centre of the
room. Pushing the shelves against the walls removes the possibility.

\textbf{The two delivery bays are in opposite corners and the station is somewhere else again.}
Parcels are nine-centimetre cubes, deliberately shorter than the height of the lidar, so a carried
parcel is never mistaken for an obstacle. The consequence is that a delivered parcel is invisible to
the navigation system. Placing deliveries close together would build a pile of obstacles the robot
cannot see, so the bays are far apart, and the station is offset from both so the robot does not
park on something it has delivered.

That reasoning is sound for parking but incomplete for driving. The station and both bays sit at the
same depth in the room, so the route home from either bay runs along the line the other occupies. A
delivered parcel is below the height of the laser and therefore invisible to the planner, and the
robot pushes it along rather than avoiding it. This was observed directly and is recorded in
Chapter~\ref{ch:limitations}; it is a placement problem, not a control problem, and the fix is to
move the station off that line.

\textbf{The two parcels on each shelf are about a metre apart.} A parcel is carried on the floor
rather than lifted, so the only thing preventing a carried parcel from knocking into a stored one is
physical separation. Lifting was considered and rejected: raising a parcel into the plane of the
laser would make the robot detect an obstacle directly in front of itself and refuse to move.

\section{The Robot}
The robot is a differential-drive platform with two driven wheels, two casters, a lidar and a
gripper link at the front. Two details of the model are worth recording.

The wheels are directly under the centre of mass. In an earlier version they were behind it, which
put a large share of the weight on a nearly frictionless caster; the robot skidded and the wheel
odometry became unreliable, which matters because the localization system uses that odometry.

The gripper is not a mechanical claw. It is a joint provided by the simulator that can be attached
and detached on command, which binds a parcel rigidly to the gripper link. This is a simplification,
and it is the main source of the delivery inaccuracy discussed in Chapter~\ref{ch:limitations}: the
joint binds the parcel wherever it happens to be lying, which is often some distance off centre, and
releases it at that same offset.

\screenshotslot{Figures/shot-gazebo-warehouse.png}{0.95\textwidth}
{The warehouse in the simulator, with all six parcels on their shelves and the robot at its station.
This is the world every run starts from.}{fig:shotgazebo}

\screenshotslot{Figures/shot-rviz-costmap.png}{0.95\textwidth}
{The same warehouse as the robot perceives it: the laser scan lying on the walls of the known map,
and the inflation the planner keeps around every obstacle.}{fig:shotrviz}

\section{The Map}
The navigation system needs an occupancy map. Rather than producing one by driving the robot around
and building it from sensor readings, the map is generated analytically from the same wall and shelf
definitions that describe the world. This has three advantages: there is no drift, because the walls
are exactly where the simulator puts them; the process is perfectly repeatable; and the map frame is
identical to the world frame, which means every waypoint in the database can be written in world
coordinates and compared directly with positions read from the simulator. The corresponding
obligation is that the map must be regenerated whenever the geometry changes.

% ============================================================
\chapter{Implementation}
\label{ch:implementation}
% ============================================================

\section{Overview}
This chapter describes how the components are implemented, following the layers from the bottom
upwards. Three sequence diagrams show how they cooperate, one scenario each, because a single
diagram carrying every branch is harder to read than three that each follow one path to its end.
Figure~\ref{fig:seqsuccess} is the normal case; Figure~\ref{fig:seqrefused} is an order the system
declines to accept; Figure~\ref{fig:seqestop} is a running order halted by the operator. Together
they cover the three outcomes an order can have: it is delivered, it is never started, or it is
stopped.

\uploadedfigure{Figures/03-sequence-order-success.png}{0.98\textwidth}
{Scenario 1: a controller places an order and it is delivered. The decision is taken before
anything is stored, and the outcome is decided by reading the parcel's real position rather than by
the fact that every command returned success. The figures are those of order 1 of the recorded
acceptance run.}{fig:seqsuccess}

\uploadedfigure{Figures/11-sequence-order-refused.png}{0.92\textwidth}
{Scenario 2: the same request is refused because the robot has not the battery to finish it and
return. The classifier is consulted, but the policy gates decide; nothing is written to the orders
table and the robot is never sent, so a refusal costs one preview.}{fig:seqrefused}

\uploadedfigure{Figures/12-sequence-emergency-stop.png}{0.95\textwidth}
{Scenario 3: the operator halts a running order. The task manager carries \emph{two overlapping
execution occurrences}: the outer one is its worker thread, blocked waiting on a navigation result
at exactly the moment the button is pressed, and the inner one is the subscription callback, which
is free to act while the worker cannot. That is why the stop is handled on the callback rather than
on the worker. Nothing is written by the stop itself, so it does not survive a logout and the next
order is an ordinary one.}{fig:seqestop}

\section{The Database Layer}
The database layer is a single class that owns the schema and exposes domain operations. On first
use it creates the tables and inserts the seed data describing the warehouse and the catalogue, so
a fresh installation is immediately usable.

The one implementation detail that caused a real bug is worth recording. When an order is created,
the row must be committed before it is read back, because the read opens its own connection and an
uncommitted write is not visible to it. The insert and the read-back are therefore separated
deliberately.

\section{The Admission Workflow}
The admission workflow is implemented as a module of plain functions that take dictionaries and
return a decision. It has no knowledge of HTTP, of the database or of the robot. Its structure is
described in detail in Chapter~\ref{ch:complex}.

\section{The Task Manager}
The task manager executes one order at a time. It takes the oldest pending order, drives to the
pick point, attempts the grab up to three times, carries the parcel to the delivery bay, releases
it, and verifies the result before reporting the outcome. If the grab fails three times the order
fails and the robot returns to its station.

The execution runs on a worker thread which blocks while waiting for navigation results, with the
message callbacks serviced by a separate executor. This is much easier to follow than a state
machine driven entirely by callbacks, which was the alternative considered.

Two behaviours are worth describing. When the robot has no work, a timer sends it back to its
station, where it reports itself as charging and the condition model begins to recover the battery
and cool the motors. And when the system starts, the task manager repeatedly publishes the robot's
initial position until the localization system confirms it by publishing a position of its own.
This handshake replaced a fire-and-forget approach that failed intermittently: early in start-up the
localization system silently discards a position it cannot yet transform, and if all the attempts
happened to fall inside that window the robot never localized and simply sat still.

\section{The Gripper}
The gripper node exposes two services, grab and release, and translates them into commands to the
simulator's detachable joint. Its important function is not commanding but verifying, and the
verification is described in Section~\ref{sec:c3}.

\section{The Condition Monitor}
The condition monitor integrates the update rules of the condition model against what the robot is
really doing. It takes distance from the wheel odometry rather than from the commanded velocity,
which is a deliberate choice: if the robot is stuck against an obstacle the wheels turn but the
robot does not move, and a model driven by commands would report energy being spent on travel that
never happened.

It publishes the resulting condition once per second and writes the same values to a file, one file
per run. That file is the input to the machine-learning pipeline.

\section{The Prediction Service}
The prediction service loads the trained model and exposes one operation. If no model is present it
does not fail: it starts normally and reports that no model is loaded, which allows a caller to
distinguish ``the model says no'' from ``there is no model''. Those are very different situations
and confusing them would hide a deployment error.

\section{The Web Application}
The web application is the control panel for the system. It serves the order form, the preview, the
order history, the robot condition, a confirmation page for sending the robot home, a login page,
and three administrator pages: product maintenance, delivery-point maintenance and a read-only
database browser. All are rendered on the server from templates. The two pages that display live
values refresh themselves every two seconds using a standard HTML directive rather than JavaScript,
which keeps the no-JavaScript constraint intact.

The pages show data and controls only. No explanatory text appears on any of them: guidance belongs
in the documentation, where it can be read once, rather than on a panel an operator looks at many
times a day. The visual design is deliberately plain for the same reason, and colour is used only
where it carries meaning.

\subsection*{Access control and roles}
Every page requires a session, and an anonymous visitor is redirected to the login page with their
destination remembered. There are two roles. A \emph{controller} may order, track, inspect the robot,
send it home and stop it. An \emph{administrator} may additionally maintain the catalogue and the
layout, manage accounts and browse the database. The distinction is enforced on the request, not
merely hidden in the navigation, so posting directly to an administrator endpoint as a controller is
refused rather than obeyed.

Passwords are stored as salted hashes using only the standard library, and the cookie holds nothing
but a random session token, so there is no signing key that could leak. Sessions are cleared on every
launch; accounts are not.

\screenshotslot{Figures/shot-login.png}{0.72\textwidth}
{The login page, reached by requesting any page without a session. The destination is remembered, so
signing in continues to wherever the visitor was going.}{fig:shotlogin}

\subsection*{The emergency stop}
A single button on every page abandons all work and halts the robot. It has no confirmation step,
because a stop that asks whether you are sure is not an emergency stop, and it does not latch,
because a stop that stays engaged after the emergency has passed becomes an obstacle of its own. It
cancels the navigation goal in flight, brakes, fails the running order and fails everything queued
behind it. The robot is then immediately available for new work.

The button sits alone at the far right of the header with a wide margin of empty space beside it.
That spacing is the safeguard against pressing it by accident, chosen deliberately in place of a
confirmation dialogue that would slow down a real emergency.

\screenshotslot{Figures/shot-estop.png}{0.95\textwidth}
{The order list immediately after the stop was pressed. The order that was running has been failed
with the reason, and the work queued behind it has been failed too rather than left to start on its
own once the robot was released.}{fig:shotestop}

\subsection*{Sending the robot home}
Returning to the station is offered as a task rather than as a command: it is placed on the same
queue as deliveries and waits its turn. A separate immediate channel would let an operator interrupt
a delivery that was already carrying a parcel, which is precisely the situation the release logic
exists to avoid. The trip is costed and reserved like any other work, because the drive really does
spend charge, but unlike a delivery it is never refused: admission control exists to stop the robot
accepting work it cannot finish, and going home is the one job that leaves a struggling robot better
off. Only one such trip may be outstanding at a time.

\subsection*{Administration and the database browser}
The administrator pages create, update and delete products and delivery points through the same
database methods the rest of the system uses. The browser displays every table read-only, a page at a
time, with password hashes, salts and session tokens redacted.

Table names cannot be supplied as parameters in a query the way values can, so the name of the table
being viewed is checked against the database's own catalogue of tables before it is used. That check
is the defence against a crafted table name, and it is one of the behaviours pinned by the automated
tests.

\screenshotslot{Figures/shot-admin-products.png}{0.95\textwidth}
{Product maintenance. The catalogue is the record the rest of the system reads, so this page is the
only place the warehouse contents are defined.}{fig:shotadminproducts}

\screenshotslot{Figures/shot-admin-db.png}{0.95\textwidth}
{The read-only database browser, showing the user table. Password hashes, salts and session tokens
are replaced by a marker rather than displayed.}{fig:shotadmindb}

The preview page is the centre of the interface, because it is where the decision is shown to the
user before anything is committed. It displays the product, the route, the predicted energy and
time, the battery level the robot would be left with, the peak temperature the job would produce,
and the decision with its reason.

\screenshotslot{Figures/shot-order-form.png}{0.95\textwidth}
{The order form. The client chooses a product and a destination; the catalogue below shows the
weight and storage location of each product, and no coordinates appear anywhere.}{fig:shotform}

\screenshotslot{Figures/shot-preview-accepted.png}{0.95\textwidth}
{An accepted order preview. The predicted route, energy, time, resulting battery level and peak
motor temperature are shown before the order is committed, and the header reports the robot's
condition on every page.}{fig:shotaccepted}

\screenshotslot{Figures/shot-orders-page.png}{0.95\textwidth}
{The order history after four deliveries. Each row carries the predicted energy alongside the
measured result, which is what makes the two comparable afterwards.}{fig:shotorders}

\screenshotslot{Figures/shot-robot-page.png}{0.95\textwidth}
{The robot condition page. The robot's position is deliberately absent: this page answers whether
the robot is fit to work, not where it is.}{fig:shotrobot}

% ============================================================
\chapter{Complex Functionality in Detail}
\label{ch:complex}
% ============================================================

The course requires complex features to be described algorithmically. This chapter does that for
the three features classified as complex in Chapter~\ref{ch:categorization}.

\section{C1 - The Robot Condition Model}

\subsection{How the Model Works}
The model does not track a single number for the robot's health. It tracks three quantities that
each affect the other two, and it is that mutual influence which makes it a model rather than three
separate counters.

\textbf{Energy} is spent in proportion to the distance travelled. Carrying a parcel makes every
metre more expensive, and a heavier parcel makes it more expensive still, because the motors have
to move more mass. On top of that, the same journey costs more when the robot is in poor shape: a
hot motor is less efficient than a cool one, and a worn robot is less efficient than a new one.
Both of these appear as multipliers on the basic cost, and neither can ever make a journey cheaper
than the ideal case of a cool, new robot.

\textbf{Temperature} behaves the way anything with a heater and a cooler behaves. While the motors
are working, heat is produced, and more heat is produced when the robot is carrying a heavy load.
At the same time the robot is always losing heat to the air around it, and it loses heat faster the
hotter it is. The result is that temperature climbs steeply at first and then flattens out towards a
value that depends on how hard the robot is working and what it is carrying. When the robot stops
and docks, the heating stops and the temperature falls back to room temperature. This is why a
light product and a heavy product settle at noticeably different temperatures, and it is the reason
weight can decide whether a job is allowed.

\textbf{Condition} represents mechanical wear. Unlike the other two it only ever goes down, and it
goes down slowly. Two things wear the robot: the total amount of energy that has been pushed
through it, and time spent running hotter than it should. A robot that does gentle work for a long
time wears less than one that does heavy work in the heat.

The important part is how the three are connected. Wear makes the robot less efficient, which means
it uses more energy for the same job, which produces more heat, which in turn causes more wear. A
robot that has been worked hard therefore becomes progressively more expensive to run, which is the
behaviour a real machine shows and the reason the model was built this way rather than as three
independent numbers.

The same three quantities are what the admission workflow reads when it has to decide whether a new
order is possible, and what the robot condition page displays to the operator.

\subsection{Choice of Constants}
The constants were chosen so that the relationships are observable within a demonstration rather
than being technically present but invisible. The battery capacity gives a single delivery of the
heaviest product per charge, or about three of the lightest; a much larger battery would be more
realistic for a real warehouse robot but would mean the battery never became a constraint. Measured
across the whole catalogue from a full charge, the most expensive order leaves the robot at
thirty-five per cent and the cheapest at seventy-five, so every order in the catalogue is still
accepted from a full battery and none has been made impossible - the robot simply has to return to
its station far more often. The heating constant was chosen so that the
lightest product settles at a temperature safely below the limit while the heaviest settles above
it, which is what makes payload weight genuinely decide an outcome rather than merely appearing as
an unused input. This is a modelling decision and is stated as such.

Figure~\ref{fig:payloadenergy} shows the resulting relationship between payload and energy per
metre.

\uploadedfigure{Figures/fig-payload-energy.png}{0.78\textwidth}
{Energy per metre for each of the six catalogue weights, computed from the model.}
{fig:payloadenergy}

\subsection{The Duty Cycle}
The update rules describe how the robot's condition changes. A state machine decides what the robot
is allowed to do about it.

\uploadedfigure{Figures/04-state-duty-cycle.png}{0.88\textwidth}
{The duty-cycle state machine. Health states outrank activity: an overheated robot is in cooldown
whatever the task manager believes it is doing.}{fig:statemachine}

The important property of this machine is that the health states are derived from measurements, not
from commands. The cooldown state is entered because the measured temperature reached the limit,
and the admission workflow reads that state and refuses new work while it holds.

The four figures that follow are measurements, not predictions. They were collected from the robot
while it was running, using the log file the condition monitor writes once per second. Several such
runs were carried out during development, and in all of them the aim was the same: to use the
system to one hundred per cent of its capacity rather than to demonstrate a comfortable case. That
means ordering the heaviest products, sending the robot on the longest routes across the warehouse,
and continuing to place orders until something actually gave way. A run in which the battery stays
high and the motors stay cool proves very little; a run in which the battery falls towards the
reserve and the motors reach their limit is the one that shows whether the model and the policy
behave correctly under pressure.

The session shown here is one of those runs: every product in the catalogue delivered, heaviest
first and on the longest routes, with orders placed back to back and no idle time allowed. Over its
course the battery fell from full to twenty-eight per cent across several charge cycles, the motor
temperature climbed to sixty-seven degrees, and mechanical condition degraded measurably. Every one
of those events was produced by the system itself rather than staged.

One thing the session does \emph{not} show is worth recording, because it would be easy to present
the figures as though it did. The motors never reached the temperature limit, and the robot was
never forced into a cooldown by overheating. That is a direct consequence of the battery size. A
charge is worth roughly one delivery, so the robot must return to its station and charge between
jobs, and charging is also when it cools fastest; heat therefore cannot accumulate across
consecutive deliveries the way it could on a larger battery. The temperature gate is still present,
still correct and still tested, but on these constants the reserve gate reaches its limit first and
refuses the work before the motors ever become the binding constraint. The cooldown state does
appear in the figures, but in its milder form: a robot that is idle and still warm reports itself as
cooling rather than as ready.

\uploadedfigure{Figures/fig-battery-session.png}{0.95\textwidth}
{Battery over one session, taken from the run log. The shaded intervals are the periods when the
robot was carrying a parcel, which are visibly steeper than the unloaded intervals.}{fig:battery}

\uploadedfigure{Figures/fig-temperature-session.png}{0.95\textwidth}
{Motor temperature over the same session. The temperature crosses the limit during a sequence of
deliveries, the robot enters cooldown, and the temperature then returns to ambient once it is
docked.}{fig:temperature}

\uploadedfigure{Figures/fig-duty-states.png}{0.95\textwidth}
{The duty-cycle states over the same session.}{fig:duty}

\uploadedfigure{Figures/fig-condition-wear.png}{0.95\textwidth}
{Mechanical condition over the same session, degrading with energy throughput and with time spent
above the warning temperature.}{fig:wear}

\section{C2 - Order Admission and Cost Preview}

\subsection{The Procedure}
The admission workflow is the decision procedure of the system. Figure~\ref{fig:activity} gives it
as an activity diagram.

\uploadedfigure{Figures/05-activity-admission.png}{0.86\textwidth}
{Activity diagram of the admission workflow. The four physical gates run before the classifier is
consulted.}{fig:activity}

In pseudocode:

\begin{lstlisting}[language=Python,caption={The admission decision.},label={lst:admission}]
def decide(product, destination, station, robot_state, predictor):
    # 1. Resolve and cost the job.
    legs   = route_legs(robot_position, product.pick_point,
                        destination, station)
    energy = energy_wh(legs.order, product.weight,
                       robot_state.temperature, robot_state.condition)
    ret    = energy_wh(legs.return_leg, 0.0,
                       robot_state.temperature, robot_state.condition)
    predicted, peak = simulate_route(robot_state, legs.order, product.weight)

    # 2. Physical gates, cheapest and most certain first.
    if product.stock <= 0:
        return REFUSED, "out of stock"
    if battery_after(energy + ret) < RESERVE_PERCENT:
        return REFUSED, "could not return to the station"
    if peak >= T_MAX:
        return REFUSED, "would overheat"
    if robot_state.condition < CONDITION_MIN:
        return REFUSED, "maintenance required"

    # 3. The learned verdict, consulted only once the physics allows it.
    if predictor is not None:
        verdict = predictor(features(product, cost, robot_state))
        if verdict is not None and not verdict.feasible:
            return REFUSED, "predicted infeasible"

    return ACCEPTED, cost_summary
\end{lstlisting}

\subsection{Why the Order of the Gates Matters}
The ordering in Listing~\ref{lst:admission} is the substance of the design, not an implementation
detail.

The physical gates run first because they are certain and cheap. They are also the gates whose
\screenshotslot{Figures/shot-stock-consumed.png}{0.95\textwidth}
{The order page after a delivery. The delivered product remains visible in the catalogue at zero
stock, marked as delivered, so its disappearance from the order form is explained rather than
mysterious.}{fig:shotstock}

refusals a human can act on: ``out of stock'' and ``needs charging'' tell the operator what to do,
whereas ``the model predicted this would fail'' does not. Running them first means an actionable
refusal is never masked by a statistical one.

The classifier runs last, and because it runs last it can only ever refuse an order that the
physics has already allowed. It can never approve one that the physics rejected. This is the
concrete sense in which the model is a component inside the system's logic rather than the logic
itself, and it is enforced by a test that supplies a classifier which always answers ``feasible''
and asserts that an order with a flat battery is still refused.

Figure~\ref{fig:shotrefused} shows this happening on the running system. The robot had already
completed four deliveries and its motors were at 68 degrees. Asking it for the heaviest product
again produced a refusal from the thermal gate, and the page states that the decision was taken by
the safety policy with no model consulted, because the gate refused the order before the classifier
was ever reached.

\screenshotslot{Figures/shot-preview-refused.png}{0.95\textwidth}
{A refused order. The predicted peak temperature exceeds the limit, so the thermal gate refuses
before the classifier is consulted at all - which is exactly the ordering the workflow is designed
around.}{fig:shotrefused}

\subsection{Work Already Promised}
A decision taken against the battery the robot reports is a decision about a robot standing still.
If two orders have already been accepted and are waiting their turn, the charge they will consume is
spoken for, and judging a third order against the present reading would accept work the robot cannot
actually reach.

The workflow therefore projects the robot forward before any gate runs. The energy every unfinished
order is expected to consume is summed, the hottest peak among them is carried across, and the gates
judge the new order against what would be left afterwards rather than against what is showing now.
A refusal caused this way says so, naming the number of jobs already in progress, and the preview
adds a line reporting the charge the order would actually start from - without it the figure
contradicts the battery shown elsewhere on the page and looks like a fault.

What makes this trustworthy is where the reservation lives. It is not stored as a running total
alongside the orders; it is derived from the orders table each time it is needed, by summing the ones
that have not reached a terminal state. An order that completes, fails or is refused therefore
releases its reservation as a consequence of its status changing, with nothing to remember to undo.
A stored total would have to be decremented by every path that ends an order, and the first path
that forgot would leave the robot permanently pessimistic about itself.

\screenshotslot{Figures/shot-preview-reserved.png}{0.95\textwidth}
{A preview taken while other work is already in progress. The battery the order would start from is
reported separately from the battery the robot is showing now, and the refusal names the jobs whose
charge is already spoken for.}{fig:shotreserved}

The projection happens before the first gate in Figure~\ref{fig:activity}, so every gate in that
diagram judges the same robot.

\subsection{Behaviour When the Model Is Unavailable}
If the prediction service cannot be reached, the client returns nothing and the workflow proceeds
to accept the order on the strength of the physical gates alone, recording that the decision was
taken by policy rather than by the model. The system therefore has two operating modes, one with
the model and one without, and both produce decisions.

\section{C3 - Physical Delivery Verification}
\label{sec:c3}
The first two complex functionalities decide what the robot should do. This one decides what the
system is allowed to \emph{believe} it has done, and it is the reason no false success has ever
reached the database.

\subsection{The Problem It Solves}
A delivery robot is a distributed system whose components report on themselves. The navigation
action returns \texttt{SUCCEEDED} when it thinks it has arrived. The gripper service returns success
when it has published an attach message. The task manager marks an order complete when its own
steps have run. Each of those is a statement about a \emph{command}, not about the warehouse.

Early in the project the system reported three deliveries out of three while no parcel had moved at
all. Every command in the chain had returned success and every status field agreed with every
other. The fault was invisible precisely because the only evidence being consulted was produced by
the same pipeline that had failed. Nothing in a system built this way detects that class of failure
from the inside.

C3 is the answer: before an order may be called complete, the outcome is checked against a source
of truth that the command pipeline does not produce and cannot influence --- the physics engine
itself.

\subsection{The Algorithm}
Verification runs at two levels, and both must pass. Figure~\ref{fig:seqverify} gives the full
sequence; the steps are as follows.

\textbf{Level 1 --- did the gripper really let go?}
\begin{enumerate}[nosep]
  \item Publish a detach message for the parcel, then wait a settling interval.
  \item Read the state the \emph{joint itself} reports. The simulator's detachable joint publishes
        the literal strings \texttt{attached} and \texttt{detached} on its own topic; this is the
        authority, not the fact that a detach message was sent. That topic carries a message only
        when the state \emph{changes}, so the subscription is created when the node starts rather
        than when a release is requested: a subscriber created at the moment of the release would
        miss the very transition it exists to observe.
  \item If the joint still reports \texttt{attached}, repeat from step 1, up to three attempts.
  \item If it still will not let go, the order fails, the queued orders are failed with a reason
        rather than left pending, and a recovery routine keeps retrying so that the robot heals
        itself as soon as the joint reports honestly again.
  \item If the joint has never reported at all, the release is treated as successful but the
        recorded evidence \emph{says that it is unverified}, rather than presenting a guess as a
        confirmation.
\end{enumerate}

\textbf{Level 2 --- is the parcel actually at the bay?}
\begin{enumerate}[nosep, start=6]
  \item Wait one second for the physics to settle, so that a parcel still in motion is not measured
        mid-fall.
  \item Read the parcel's position from the simulator.
  \item If no position has ever been received, fail the order with ``cannot confirm delivery''.
        Absence of evidence is treated as failure, not as permission to proceed --- this is the
        exact gap the original fault escaped through.
  \item Compute the distance between the parcel and the delivery bay.
  \item If that distance exceeds the tolerance of 0.7 m, fail the order with the measured distance
        in the message and send the robot home. The order is failed even though every command in it
        succeeded.
  \item Only if the distance is within tolerance is the order marked complete, and the measured
        distance is recorded with it.
\end{enumerate}

\uploadedfigure{Figures/10-sequence-verification.png}{0.95\textwidth}
{C3, physical delivery verification. Two independent sources of truth are consulted --- the joint's
own reported state and the parcel's position in the physics engine --- and neither is produced by
the command pipeline that asked for the work. Every path that does not end in confirmed evidence
ends in a failed order.}{fig:seqverify}

\subsection{Why This Is Complex Rather Than a Status Check}
A status check asks a component whether it succeeded. This asks a different question of a different
system, and then decides which answer wins.

Three properties make it more than a comparison. First, it is a \emph{cross-validation between two
independent sources of truth}: the command pipeline and the physics engine are separate systems,
and the design consists of deciding which one is authoritative when they disagree. Second, the
\emph{failure semantics are inverted} relative to ordinary error handling: the default outcome is
failure, and success must be earned by positive evidence, which is why a missing measurement fails
an order rather than passing it. Third, it required identifying what counts as admissible evidence
at all --- the project moved from proximity, which passed trivially when nothing happened, to the
joint's own reported state, which cannot.

That last point is worth stating plainly, because it is where the engineering lies. An earlier
version of the grab check measured the gap between parcel and robot, published an attach, and
measured again, accepting the grab if the gap was unchanged. If the attach silently did nothing,
neither body moved, the two measurements were identical, and the check passed. The check could not
fail. Replacing it required recognising that a measurement taken from the wrong side of the failure
tells you nothing, and finding a signal that is generated by the mechanism being tested rather than
by the code testing it.

\subsection{What It Costs and What It Buys}
The cost is real: a delivery that physically succeeded can be failed by this check if the parcel
lands outside the tolerance, and the robot is sent home with the order marked failed. That is
deliberate. The alternative is a system whose records and whose warehouse disagree, and the project
takes the position that a false success is worse than an honest failure.

What it buys is that every completed order in the database is backed by a measurement taken from
the simulator, and the acceptance suite can therefore assert physical outcomes rather than status
fields. It is also what makes the reliability figures in Chapter~\ref{ch:testing} meaningful: a
delivery counted as successful there is one where the parcel was measured at the bay.

% ============================================================
\chapter{The Machine Learning Component}
\label{ch:ml}
% ============================================================

\section{Purpose and Scope}
The machine-learning component answers a single question: given the robot's condition and the size
of a proposed job, will the job complete within the battery and thermal limits? It is a binary
classifier, and it is deliberately the only model in the system. An earlier design considered three
models, one for energy, one for peak temperature and one for remaining runtime, but a single
feasibility classifier covers the decision the system actually has to take, and one model is far
easier to explain, test and justify than three.

\section{Features}
The model receives five numbers, all physically meaningful.

\begin{table}[H]
\centering
\begin{tabular}{@{}lp{0.62\textwidth}@{}}
\toprule
\textbf{Feature} & \textbf{Why it is relevant} \\
\midrule
Battery percentage & Determines whether the robot can finish and still return. \\
Motor temperature & Determines how much thermal headroom remains. \\
Mechanical condition & A worn robot consumes more energy for the same work. \\
Payload mass & Heavier products consume more energy and generate more heat. \\
Route distance & The total travel including the return leg. \\
\bottomrule
\end{tabular}
\caption{The classifier's features.}
\label{tab:features}
\end{table}

\section{Training Data}
Training a model requires thousands of labelled examples, and a single delivery in the simulator
takes about a minute. Collecting enough real deliveries would therefore have taken days of
continuous running, which was not practical within the project.

The training set is instead produced by a generator that creates \emph{randomly generated data
intended to simulate how the system would be used in the real world}. Each generated example is one
hypothetical order: a robot in some starting condition, carrying one of the six catalogue products,
over a route of a plausible length. The random values are not arbitrary. The battery, temperature
and wear levels are drawn from the ranges the robot genuinely passes through during operation, the
weights are exactly the six weights in the catalogue, and the route lengths are limited to the
distances this warehouse can actually produce, from a short trip to a nearby shelf up to the longest
corner-to-corner journey. Generating a route of fifty metres would teach the model about a journey
the robot can never be asked to make.

What makes the generated data meaningful rather than invented is that it is produced by
\emph{the same condition model the live robot runs}: the generator imports that module instead of
reimplementing the behaviour. A generated run is therefore what the real robot would have done in
that situation, according to the same rules it uses in operation. The output is written in the same
file format the robot produces during a real run, so generated and measured runs are interchangeable
and could be combined into a single training set as more real data becomes available.

\uploadedfigure{Figures/fig-ml-pipeline.png}{0.95\textwidth}
{The machine-learning pipeline. The generator and the trainer are disposable; the running system
depends only on the exported model.}{fig:mlpipeline}

Each generated example samples a plausible starting condition, a payload drawn from the six
catalogue weights and a route length within the range this warehouse can produce, then runs the
condition model forward and labels the example according to whether the job would have finished
within the limits. A small amount of noise is applied to the thresholds so the boundary between the
two classes is not a perfectly sharp line, which would otherwise let the classifier learn a trivial
rule.

\uploadedfigure{Figures/fig-ml-class-balance.png}{0.62\textwidth}
{Class balance of the generated training set. The set holds 4000 hypothetical orders, each
labelled by running the deterministic model forward to see whether the robot could in fact
have completed it.}{fig:mlbalance}

\section{Avoiding Information Leakage}
The generated file contains columns that the model is deliberately \emph{not} given, in particular
the battery level after the job and the peak temperature reached. Those values are produced by the
same formula that produces the label, so including them would allow the model to read the answer off
its own input and score almost perfectly while learning nothing. The model receives only quantities
that are knowable before the job runs, which is exactly the information the admission workflow has
at the moment it must decide.

\section{Model and Results}
The classifier is a random forest inside a pipeline with a feature scaler, so the exported artefact
is self-contained and the service can hand it raw physical units. On a held-out quarter of the data
the model reached an accuracy of ninety-seven per cent, and five-fold cross-validation on the
training split gave an F-measure of 0.945.

The training set is not evenly split between the two classes, and deliberately so: with the battery
sized as it is, roughly a third of the generated orders are feasible and the rest are not. A
classifier trained on a set where almost everything succeeds learns to say yes, which is worth
nothing to a workflow whose entire purpose is to be able to say no.

\uploadedfigure{Figures/fig-ml-confusion.png}{0.60\textwidth}
{Confusion matrix on the held-out test split, the 1000 runs of the 4000 not used for training.}{fig:confusion}

\uploadedfigure{Figures/fig-ml-importance.png}{0.82\textwidth}
{Relative importance of the five features.}{fig:importance}

Figure~\ref{fig:importance} shows that the battery level dominates, followed by the length of the
route. This is consistent with the physics: the reserve constraint is the tightest of the four, and
route length is what converts a healthy battery into an insufficient one. Payload and temperature
matter less directly, influencing the outcome largely through their effect on energy consumption.

\section{An Honest Statement of What the Model Means}
The accuracy figures above describe how well the model reproduces the behaviour of the generator
that produced its training data. They do not describe how well it predicts the behaviour of a real
robot, because it was never trained on one. This is stated plainly because presenting the figure as
evidence of real-world predictive power would be misleading.

What the component demonstrates is the integration of a learned model into a decision workflow: the
construction of features from live system state, the consultation of a model as one input among
several, a policy able to overrule it, and a defined behaviour when it is unavailable. Should real
run data become plentiful, the file format is already identical and the model could be retrained on
it without any change to the rest of the system.

% ============================================================
\chapter{Testing and Validation}
\label{ch:testing}
% ============================================================

\section{Strategy}
Testing is organised in three layers, in increasing order of realism and decreasing order of speed.

\uploadedfigure{Figures/fig-test-pyramid.png}{0.68\textwidth}
{The three layers of the test strategy.}{fig:pyramid}

\textbf{Unit tests} exercise pure logic: the condition model, the admission policy and the database.
They require no simulator and no network.

\textbf{Integration tests} exercise the joins between components: an HTTP request producing a
database row and a message to the robot, with the product resolved to its simulator model and its
pick and delivery coordinates so the robot side never needs the catalogue; and the login gate
applied to every page that requires a session. The robotic middleware and the prediction service
are replaced by fakes, so the tests remain fast and deterministic.

\textbf{Acceptance tests} run against the complete live system with the simulator running, and
verify physical claims against the simulator rather than against status fields.

\section{Automated Test Suites}
The automated suites contain 16 tests and run in a few seconds. One of them - the login gate - is
applied to each of the eight pages that require a session, so the 16 tests execute as 23 cases.

The suite is deliberately small. It is a representative set rather than an exhaustive one: one test
per behaviour the system genuinely depends on, each of which can be stated in a sentence and
defended. A larger suite catches more regressions, and that trade is made consciously here in favour
of a body of tests that is small enough to be understood in full. Physical behaviour is not covered
by these at all; it is covered by the acceptance suite, which runs against the live simulator and is
the larger of the two.

Three of the sixteen are integration tests and need a word of explanation, because they are the only
part of the suite that is not a plain function call. The application builds its database connection,
its link to the robot and its client for the prediction service at the moment it is imported. A test
that wants a throwaway database and no robot must therefore put those substitutes in place and then
re-import the application, which is what the fixture does: it points the application at a temporary
database file, replaces the robot link with a stand-in that records what it was asked to publish
instead of sending it, and reloads the module so those choices take effect. The test then drives the
application over HTTP exactly as a browser would.

\uploadedfigure{Figures/fig-test-composition.png}{0.88\textwidth}
{Composition of the automated test suite.}{fig:testcomp}

The tests are written to pin the \emph{relationships} the system depends on, not only its
arithmetic. Examples include: a heavier payload must cost more energy and run hotter; a hot or worn
robot must make the same job more expensive; simulating a hypothetical order must not modify the
real robot's condition; and the classifier must not be able to approve an order that the physical
gates rejected. If any of those relationships were inverted, the system would still run and would
still produce plausible-looking decisions, which is exactly why they are asserted explicitly.

\section{Acceptance Testing}
The acceptance suite starts from a freshly launched system, submits real orders and checks the
outcome against the simulator. It refuses to run if the parcels are not in their starting positions,
because a world left over from a previous run would cause failures for the wrong reason.

\uploadedfigure{Figures/fig-acceptance-results.png}{0.95\textwidth}
{Results of the acceptance suite against the live system.}{fig:acceptance}

The same verification is applied to a single order: the admission verdict and its reasoning are
recorded, then each state the order passes through, and finally the parcel's real position is read
back from the simulator and compared with the delivery bay. A delivery is never reported as
successful on the strength of a status field alone.

% The data rows are generated from the recorded run by
% tools/report/make_results_table.py, so this table and the figure above it can
% never describe two different runs. Regenerate both with make_figures.py.
\begin{longtable}{@{}p{0.09\textwidth}p{0.38\textwidth}p{0.47\textwidth}@{}}
\toprule
\textbf{Req.} & \textbf{Check} & \textbf{Measured result} \\
\midrule
\endhead
% Generated by tools/report/make_results_table.py - do not edit by hand.
% Regenerate after any acceptance run:
%     ./robofetch_venv/bin/python tools/report/make_results_table.py
% Data rows only - the surrounding longtable lives in the report itself.
SETUP & API up, robot connected, all 6 parcels on their shelves & database=robofetch.db, AI=http://localhost:8001 \\
UC1 & Catalogue and layout are served from the database & 6 products, 2 destinations, weights [0.3, 0.6, 1.1, 1.8, 3.2, 4.5] \\
UC12 & An anonymous visitor is sent to the login page & GET /robot while signed out -> 303 (expected 303 redirect) \\
UC12 & A controller can drive the robot but not administer it & controller: /robot -> 200, /admin -> 403 (expected 200, 403) \\
UC12 & An administrator reaches the admin section and the database & admin: /admin -> 200, /admin/db?table=products -> 200 \\
UC13 & An administrator can create, update and delete a product & created=True, updated=Renamed widget @ 2.6 kg, deleted=True \\
UC9 & Order preview returns a cost and a verdict & 8.6 m, 6.11 Wh, peak 56.8 C -> accepted (by model) \\
NFR1 & Preview (including the ML call) responds within 2 s & 416 ms \\
C2 & The classifier is consulted inside the admission workflow & prediction service at http://localhost:8001 reachable=True, decided\_by=model \\
NFR3 & All pages render server-side with no JavaScript & 8 pages all returned 200 and none contain a script tag \\
UC3 & An order can be submitted by product ID & order 1 for SKU-1001, accepted by model, predicted 3.71 Wh \\
C2 & Work in flight is reserved against the next admission decision & robot reports 100.0\%, next order costed from 56.7\% with 1 order(s) holding 4.76 Wh \\
UC4 & An order can be tracked to a terminal state & order 1 -> completed: delivered (0.49 m from target) \\
C2 & The reservation is released when the order finishes & 0 order(s) still reserved after order 1 ended \\
UC5 & The parcel is physically delivered (Gazebo ground truth) & parcel\_1 at (-2.61, -1.94), 0.47 m from the bay, moved 3.31 m \\
FR16 & A delivered product is no longer orderable & 5 of 6 products orderable; re-ordering SKU-1001 -> refused ('SKU-1001 is out of stock') \\
FR10 & Battery, temperature and condition are tracked & battery 73.7\%, temp 52.0 C, condition 99.97\%, state=idle \\
FR12 & The run is logged to CSV for ML training & run\_20260830\_164747.csv: 36 samples \\
FR5 & An order the robot cannot serve is refused, with a reason & status=refused, reason='SKU-1002 is out of stock' \\
FR11 & An idle robot returns to its station and recharges & idle at 73.7\% -> docked and charging, 100.0\% -> 100.0\% \\
FR6 & Orders are served oldest-first (FIFO) & submitted [1, 2], started [1, 2] - the earlier order ran first \\
FR8 & A failing grab is retried 3 times, then the order fails & order 3 -> failed after 3 attempts: could not grab parcel\_6 in 3 attempts; returning to station \\
UC11 & The emergency stop halts the robot (Gazebo ground truth) & robot moved 0.000 m in the 5 s after it settled; listening=True \\
UC11 & The interrupted order is failed, not left pending & order 4 -> failed: halted by the emergency stop \\
UC11 & Nothing stays engaged - the robot accepts work immediately after & next order -> accepted: '8.6 m, 3.5 Wh, \textasciitilde{}56 s; battery 45\% after returning to the station' \\
UC6 & Analytics agree with the orders table & reported 2/5 completed and 0 refused; counted 2 and 0; rate=0.4 \\

\bottomrule
\caption{Acceptance results, all 26 checks passing against the live simulator.}
\label{tab:acceptance}
\end{longtable}

\section{Measured Results}

\subsection{Delivery Accuracy}
Every delivery was measured by reading the parcel's position from the simulator and comparing it
with the bay.

\uploadedfigure{Figures/fig-delivery-accuracy.png}{0.88\textwidth}
{Measured delivery distances against the acceptance tolerance.}{fig:accuracy}

The measurements cluster between 0.18 and 0.49 metres, with a mean of 0.35, against a tolerance of
0.7. That tolerance was itself tightened from 1.1 metres once the accuracy improved, because the
old one no longer discriminated between a good delivery and a poor one.

Earlier in the project these same measurements sat between 0.59 and 0.77 metres, and the cause
turned out not to be where it appeared to be. The fault was at the moment of the grab rather than
the moment of release: the robot arrived at the pick point without facing the parcel, the gripper
accepted a parcel lying anywhere within reach in any direction, and the joint then carried and
released it at whatever offset it had been seized with. The delivery inherited the error made at
the shelf. Approaching the pick point facing the shelf, refusing a grab outside a narrow arc, and
stopping one carry-offset short of the bay together removed most of it.

\subsection{Prediction Quality}
Because the admission estimate is stored on each order, the predicted route length can be compared
with the distance actually travelled.

\uploadedfigure{Figures/fig-predicted-vs-actual.png}{0.56\textwidth}
{Predicted against measured route length.}{fig:predvsactual}

The predictions are consistently slightly longer than the measured distances, which is expected: the
estimate is a straight-line calculation between waypoints, and the planned path is close to but not
exactly straight.

\subsection{Response Time}
The order preview, which includes resolving the product, computing the route, simulating the
condition model forward and calling the prediction service over HTTP, was measured at 416
milliseconds in the recorded acceptance run, against a requirement of two seconds. The figure
varies by a few tens of milliseconds between runs, since it includes an HTTP round trip on a
machine that is also running a physics simulation; what the requirement asks is that it stay far
below two seconds, and it does.

That figure was originally much worse, and the reason is described in
Section~\ref{sec:coldstart}: the very first preview after a restart took 1551 milliseconds and the
prediction timed out entirely. Loading the model when the service starts rather than when the first
request arrives brought the measurement down to its present value, and the check is now part of the
acceptance suite so a regression would be caught rather than noticed.

\subsection{A Requirement That Was Not Met}
\label{sec:nfr4}
Writing an acceptance criterion for every non-functional requirement, rather than stating the
requirement alone, immediately produced a result that the previous version of this report did not
contain: NFR4 was not being satisfied.

NFR4 asked that telemetry be recorded at least once per second while the robot is active. Measured
across the recorded run, the median interval between consecutive samples is 3.22 seconds, the worst
is 19.39 seconds, and every one of the 266 intervals taken while the robot was working exceeds one
second. The requirement was failing on every sample, and nothing detected it, because the
requirement had never been expressed as something that could be checked.

The cause is not a missing feature. The condition monitor declares a publication period of one
second and creates its timer accordingly, but it runs on simulated time, so the timer fires once
per simulated second. The row it writes is stamped with the wall clock. The simulator advances at
roughly a third of real time under the processing load described in
Chapter~\ref{ch:environment}, and three point two is that ratio. The component behaves exactly as
specified; the requirement and the timestamp disagreed about which clock they meant, and nothing in
the original wording forced that disagreement into the open.

Two conclusions follow, and the second is the more useful.

First, the requirement is now stated on a named clock, and the acceptance criterion says which one.
Measuring against simulated time makes it satisfied and honest, because the robot's own model
integrates in that frame.

Second, and more generally, this is the clearest evidence in the project for why an unmeasurable
requirement is not a requirement at all. NFR4 read as a firm commitment for the whole of
development and was quietly false throughout. The five requirements around it were being checked
continuously by the acceptance suite; this one was not, purely because nobody had written down how
it would be checked. The discipline applied to the functional requirements, where every one carries
an observable outcome, is what the non-functional requirements had been missing.

% ============================================================
\chapter{Development Environment and Engineering Challenges}
\label{ch:environment-challenges}
% ============================================================

\section{Platform and Toolchain}
The entire system, from the physics simulation to the web server, runs on a single machine. There
is no cloud component and no second computer.

The development and execution environment is Windows~11 running the Windows Subsystem for Linux
(WSL2), inside which the operating system is Ubuntu~24.04. On top of that sit the robotics
components: \textbf{ROS~2 Jazzy} as the middleware that carries messages between the parts of the
robot, \textbf{Gazebo Harmonic} as the physics simulator, and \textbf{Nav2} as the navigation
stack that handles localization, path planning and motion control.

\begin{table}[H]
\centering
\begin{tabular}{@{}lp{0.58\textwidth}@{}}
\toprule
\textbf{Component} & \textbf{Role in the project} \\
\midrule
Windows 11 with WSL2 & Host environment. The web interface is opened in a browser on the Windows
side while everything else runs inside Linux. \\
Ubuntu 24.04 & The operating system the whole system runs on. \\
ROS 2 Jazzy & Middleware. Carries orders, status updates and telemetry between the web tier and the
robot, and provides the service mechanism used to command the gripper. \\
Gazebo Harmonic & Physics simulation: the warehouse, the robot, the lidar and the attachable joint
used as a gripper. \\
Nav2 & Localization, global planning, local control and the recovery behaviours. \\
Python virtual environment & Holds the web framework and the machine-learning libraries. It is
created with access to the system packages so that the robotic middleware and the web framework can
both be imported by the same process. \\
\bottomrule
\end{tabular}
\caption{The platform the system was built and tested on.}
\label{tab:platform}
\end{table}

One detail of this arrangement is worth explaining because it influenced a design decision. Under
WSL2, graphical applications are shown through a compatibility layer, and that layer proved
unreliable for the simulator's own window. Serving the user interface as an ordinary web page,
opened in a browser on the Windows side, avoids the problem completely: the interface does not
depend on the Linux graphical stack at all. What began as a workaround turned out to be the better
architecture, because it also means the interface would work from another machine on the network
without any change.

\section{How the System Was Built and Tested}
The system consists of eight separate processes that have to start in the right order and find one
another. Doing that by hand is error-prone, so the project provides a single command that stops
anything left from a previous run, builds the workspace, and launches everything: the simulator, the
navigation stack, the robot nodes, the web service and the prediction service.

\screenshotslot{Figures/shot-launch.png}{0.95\textwidth}
{The launch console at the moment the system becomes usable. The line reporting that the robot is
localized and ready is the trustworthy signal: orders submitted before it appears are accepted and
never execute.}{fig:shotlaunch}

Testing was carried out in three situations, and each one was necessary.

\textbf{Without the simulator.} The unit and integration tests run against pure logic with the
robot and the prediction service replaced by fakes. This is where most defects were found, simply
because these tests run in seconds and could be repeated after every change.

\textbf{With the simulator, headless.} The simulator can run without drawing anything, which uses
far less processing power. Most acceptance runs were performed this way because it is faster and
more reliable.

\textbf{With the simulator and its graphical interface.} Some behaviour only appears when the
machine is under heavier load, so the final verification was repeated with the simulator's window
and the visualisation tool running. This is not redundant: the same test can pass headless and fail
with the interface open, because the extra processing demand changes the timing of everything else.

Throughout, the rule described in Chapter~\ref{ch:quality} was applied: a delivery was only
believed once the parcel's position had been read back from the simulator.

\section{Challenges Encountered}
This section describes the problems that took the most time to solve and what was learned from
each. They are included because they shaped the design far more than the features that worked
first time.

\subsection{The Robot Could Not Tell Where It Was}
The first version of the warehouse was a square room. In the visualisation the laser readings would
line up perfectly with the walls, and then the robot's estimate of its own position would suddenly
jump somewhere else, after which it drove confidently into a wall.

The cause is a property of square rooms rather than a defect in the code. Rotated by ninety degrees,
a square room produces an identical laser reading, so the localization system genuinely cannot tell
the two orientations apart, and it will sometimes settle on the wrong one. The solution was to
redesign the room: make it rectangular, which removes the rotational ambiguity, and give the shelves
different sizes and asymmetric positions so that every part of the room looks distinctive to the
laser. The shelves therefore serve two purposes at once, as storage and as landmarks. Position error
fell from metres, with sudden jumps, to a few centimetres with no jumps at all.

The lesson was that some problems that present themselves as software bugs are really problems in
the environment the software is asked to work in.

\subsection{The Robot Trapped Itself}
A gap was left between two shelves and the wall behind them. The robot drove into it, could not turn
around, and became wedged. Its position estimate then drifted by nearly three metres, and once that
had happened the planner refused every goal it was given, including one in the empty middle of the
room, because as far as it could tell the robot was surrounded.

The robot is forty-four centimetres wide and the planner keeps a safety margin around obstacles, so
any corridor narrower than roughly 1.2 metres is a trap rather than a route. The fix was to push
every shelf flush against a wall so that no such corridor exists. This is now a rule the layout must
respect, and it is the reason the shelves in Figure~\ref{fig:warehouse} touch the walls.

\subsection{The System Reported Deliveries That Never Happened}
The most instructive failure of the project. The system reported that three deliveries out of three
had succeeded. No parcel had moved.

Every command in the sequence had returned success, and every status message agreed. The check that
was supposed to confirm the delivery compared the parcel's position with the robot's, and that check
passed trivially: if nothing happens, the parcel is still exactly where it was, which is close
enough to satisfy a proximity test.

The fix was to compare the parcel's position against the \emph{delivery bay} rather than against the
robot, reading the position from the simulator itself. From that point on the project treated the
software's own reports as claims and the state of the physics engine as evidence. This single
incident is responsible for the verification approach used throughout the rest of the system.

\subsection{A Test That Could Not Fail}
A closely related problem appeared one level lower, in the gripper. Its check measured the distance
between the parcel and the robot before commanding the attachment and again afterwards, and accepted
the grab if the parcel was still close.

If the attachment silently did nothing, neither the robot nor the parcel moved, so the second
measurement was identical to the first and the check passed. The gripper reported a successful grab
and the robot drove away empty. The fix was to stop inferring and start asking: the simulator's
attachable joint reports its own state, either attached or detached, and that report is now the
authority.

A check that cannot produce a negative result is worse than having no check, because it manufactures
confidence without providing evidence.

\subsection{The Gripper Stopped Responding Entirely}
At one point the gripper stopped answering requests altogether. The process was alive and using a
full processor core, but it responded to nothing, not even to a simple query of its own settings.

Inspecting the process showed that its main thread was busy in Python while every other thread was
waiting for it. The cause was message volume: four position publishers were each sending updates
twenty times a second, and all of them were handled by the same Python event loop in that process.
Roughly seventy messages a second was more than it could process, and once it fell behind it never
caught up.

The fix was to reduce the publication rate to five times a second, which is more than enough for the
purpose those messages serve. The lesson is that in a message-based system the total message rate
into a process is a resource that has to be budgeted, and that a process can be completely
unresponsive while appearing perfectly healthy from the outside.

\subsection{Orders That Disappeared}
Submitting an order returned a normal response, and nothing happened. The order simply sat waiting
forever.

The cause was a process left running from a previous session. The script that stops the system
matched processes by name, and the web service runs under a generic interpreter name, so it was
never stopped. It kept the network port occupied, which meant the newly started web service failed
immediately, while requests to that port were still answered normally by the old process, which was
connected to a different database and to no robot at all.

Two things were changed. Shutdown now matches those services by their full command line, and
the start script refuses to launch if a port is still occupied, naming the process that holds it.
The web service also gained a health check that reports whether a robot is genuinely listening,
which turns an invisible failure into a visible one. The same problem recurred later when a second
service was added and the pattern had not been generalised, which is why the matching rule is now
written to cover services added in the future.

\subsection{The Parcel Fought the Robot}
Early deliveries frequently failed part way. The parcel is attached to the robot wherever it happens
to lie, which is often off to one side, and at ordinary friction, dragging it along the floor
produced enough resistance for the navigation system to decide the robot had stopped making progress
and to abandon the goal.

Reducing the friction of the parcels to almost nothing solved it, and delivery success improved
immediately. It is a simulation-specific compromise, and it is recorded honestly among the
limitations, because it is the reason the catalogue weight is used for the energy model rather than
being applied to the simulated object.

\subsection{An Upgrade That Broke Every Page}
After a library upgrade, every page of the web interface returned an internal error while the
stylesheet still loaded perfectly. The error message pointed deep inside the template system and
mentioned an unrelated type problem.

The cause was that the function used to render a template had changed the order of its arguments in
a newer version. The old call passed the template name first, and the new version interpreted that
name as something else entirely. Nothing about the error message suggested it.

A second, similar problem appeared at the same time: files were being served through a mechanism
that refuses to follow links pointing outside its own directory, and the build system installs files
as exactly such links. The result was that pages loaded but every stylesheet and image returned
``not found''.

Both were quick to fix once understood, and both took far longer to diagnose than to correct. The
general lesson is that an error message often describes the symptom at the point where a value
finally became unusable, not the point where the mistake was made.

\subsection{The First Prediction of Every Session Was Lost}
\label{sec:coldstart}
This one was found by the acceptance suite rather than by a person, which is the best argument for
having written it.

The check that confirms the classifier is actually consulted began failing, while every later order
in the same run reported that the model had been used. The prediction service was running, the
model was loaded, and a manual request answered instantly.

The cause was that the service loaded its model on the first request rather than at start-up.
Importing the machine-learning library and reconstructing a two-hundred-tree forest takes longer
than a second, which is longer than the calling side is willing to wait. The very first prediction
after a restart therefore timed out, the caller correctly fell back to the deterministic policy, and
every subsequent request was fast because the model was by then in memory. The system behaved
exactly as designed: it degraded instead of failing. The problem was that it degraded every single
session, for the first order, silently.

Loading the model when the service starts fixed it. The measured preview time fell from 1551
milliseconds to 283, and the acceptance suite returned to a full pass.

There are two lessons. The first is that a graceful fallback can hide a defect, because nothing
appears to be broken. The second is that lazily initialising something expensive moves the cost to
the least convenient possible moment: the first real request, when someone is waiting.

\subsection{Summary of the Challenges}

\begin{longtable}{@{}p{0.30\textwidth}p{0.28\textwidth}p{0.34\textwidth}@{}}
\toprule
\textbf{Problem} & \textbf{Cause} & \textbf{What changed} \\
\midrule
\endhead
Position estimate jumped & A square room looks identical when rotated & Rectangular room with
asymmetric shelves \\
Robot wedged and stopped & A corridor narrower than the robot can turn in & Shelves pushed flush
against the walls \\
Success reported without deliveries & The check compared the parcel with the robot, not the
destination & Verify against the simulator's own state \\
Grabs reported when nothing was held & The check could not produce a negative result & Use the
joint's own reported state \\
Gripper stopped responding & Too many messages for one process & Reduced publication rate \\
Orders accepted but never executed & A leftover process holding the port & Stop script matches full
command lines; health check reports the robot link \\
Deliveries abandoned part way & Dragging the parcel looked like a stall & Near-zero parcel
friction \\
Every page failed after an upgrade & Changed argument order in a library, and links rejected by
the file server & Updated the call; allowed links to be followed \\
First prediction of each session lost & The model was loaded on the first request, more slowly
than the caller would wait & Load the model when the service starts \\
\bottomrule
\caption{The main problems encountered and their resolutions.}
\label{tab:challenges}
\end{longtable}

% ============================================================
\chapter{System Evolution}
\label{ch:evolution}
% ============================================================

Section~\ref{sec:increments} describes what each increment delivered. This chapter records the
judgement behind the increments: why the second one was built as a replacement rather than an
addition, and what the sequence cost and returned.

\section{Why the Second Increment Replaced Rather Than Extended}
The first increment was a complete and working delivery robot. It accepted a request over the
network, navigated the warehouse, picked a parcel up, delivered it, verified the delivery against
the simulator and reported the result. As a first increment it did its job: it proved the
navigation, the manipulation and the web interface could be made to work together, and it produced
the environment and the verification discipline that everything after it inherited.

Its limitations became clear once it was examined against what a warehouse actually needs. Work was
requested by naming an item in the code rather than by choosing a product, and the information the
system held about itself was too thin to support any decision: the robot could carry out an order,
but it had no basis on which to judge whether it should. Some of the machinery it contained was also
more elaborate than the behaviour it produced justified.

The choice at that point was between extending the existing structure and replacing the parts that
were wrong. Replacement was chosen because the missing idea --- that the system should understand
its own physical state well enough to decide whether a job is possible --- reaches into the data
model, the interface and the robot's software at once. Patching that into the existing structure
would have produced something harder to explain than a clean second increment, and explicability is
a stated goal of this project.

What had proved itself was carried forward unchanged: the warehouse design, the navigation
configuration, and above all the rule that a delivery is believed only when the simulator confirms
it.

\section{What the Sequence Cost and Returned}
Two features were built, verified and then deleted, which is work that produced no shipped
functionality. Against that, the sequence returned three things that a single delivery would not
have.

It made the two withdrawals possible at all. Both were judged on observed behaviour rather than on
their design, and neither fault was visible on paper: serving orders by proximity looks reasonable
until a customer is overtaken, and a retry state machine looks prudent until its behaviour turns out
to equal that of a bounded loop.

It kept every claim testable as it was made. Because each increment ended in a system that ran, the
verification discipline described in Chapter~\ref{ch:testing} could be applied continuously rather
than assembled at the end, which is what caught a case where every command reported success and no
parcel had moved.

Finally, it absorbed the third increment without disturbance. Authentication, the emergency stop,
the return to the station and the reservation of accepted work were all added to an architecture
that was settled in the first increment and has not been re-cut since.
% ============================================================
\chapter{Limitations}
\label{ch:limitations}
% ============================================================


\textbf{Delivery accuracy is between roughly 0.2 and 0.5 metres.} The remaining error comes from
the simulated gripper: the detachable joint binds the parcel wherever it happens to be lying at the
moment of the grab and releases it at that same offset, so any residual misalignment at the shelf
is carried to the bay. The parking tolerance of the navigation system is a secondary contributor.

An earlier version of this document listed an alignment step before grabbing as the way to improve
this. That step has since been implemented - the robot now approaches the pick point facing the
shelf and the gripper refuses a parcel outside a narrow arc in front of it - and the measured error
fell from roughly 0.7 metres to roughly 0.35. What is recorded above is what remains after that
change, not what was there before it.

\textbf{The robot pushes parcels it has already delivered.} The station and both delivery bays lie
at the same depth in the room, so the drive home crosses the bay line. Parcels are shorter than the
laser height and so cannot be seen by the planner, and one was observed being pushed more than two
metres from its bay and left on the charging dock. The delivery itself is verified before the robot
moves, so no false success is recorded, but the parcel does not stay where it was put. Moving the
station off that line is the cheapest fix and is not yet done.

\textbf{Only the pages are behind the login; the interface for scripts is not.} The browser pages
all require a session, but the machine-readable endpoints do not, because the acceptance suite and
the automated checks use them and this is a single-user simulator on one machine. On a shared network
those endpoints would need the same check and every script would need a credential.

\textbf{The catalogue is the record of stock, and the simulator is not consulted.} If the catalogue
says a shelf holds four of something while the simulator holds one, ordering the missing ones simply
fails after three grab attempts. Synchronising the two was considered and rejected as more machinery
than a demonstration warrants.

\textbf{The catalogue weight is not the simulated mass of the parcel.} The parcels in the simulator
are light and have very low friction. That setting is load-bearing: with normal friction a dragged
parcel resisted enough for the navigation system's progress checker to abort the goal, and delivery
success improved dramatically when it was reduced. The catalogue weight therefore drives the energy
model and the classifier but not the physics of the simulation. This is a deliberate simplification
and is recorded as one.

\textbf{Reliability depends on available processing power.} On a heavily loaded machine, navigation
goals abort and localization degrades. This is a limit of running a physics simulation, a navigation
stack and a web application on one machine, not a defect in the application logic.

\textbf{Cancelling an order does not stop a robot that has already started.} The cancellation is
recorded in the database, but the order was handed to the robot when it was accepted, and the robot
completes it. A proper fix requires a cancellation channel to the robot.

\textbf{The map is fixed and known in advance.} The system does not build a map of an unknown
environment.

\textbf{The classifier is trained on generated data} and therefore reflects the assumptions of the
generator rather than the behaviour of a real robot.

% ============================================================
\appendix
% ============================================================

\chapter{Requirements Traceability}
\label{app:traceability}

\begin{longtable}{@{}p{0.08\textwidth}p{0.34\textwidth}p{0.50\textwidth}@{}}
\toprule
\textbf{Req.} & \textbf{Where it is satisfied} & \textbf{How it is verified} \\
\midrule
\endhead
FR1 & Product table and its access methods & Acceptance check UC1 \\
FR2 & Layout tables & Acceptance check UC1 \\
FR3 & Order endpoint and order form & Integration test; acceptance check UC3 \\
FR4 & Preview endpoint and preview page & Acceptance check UC9 \\
FR5 & Policy gates in the admission workflow & Unit tests on the stock and reserve gates;
acceptance check FR5 \\
FR6 & Oldest-first selection in the task manager & Unit test on ordering; acceptance check FR6 \\
FR7 & Task manager execution sequence & Acceptance checks UC4 and UC5 \\
FR8 & Bounded retry loop in the task manager & Acceptance check FR8 \\
FR9 & Verification step against the simulator & Acceptance check UC5 \\
FR10 & Condition monitor and telemetry tables & Unit tests on the model; acceptance check FR10 \\
FR11 & Idle timer and return-to-station behaviour & Acceptance check FR11 \\
FR12 & Per-run log writer & Acceptance check FR12 \\
FR13 & Order history page and endpoint & Acceptance check UC4 \\
FR14 & Catalogue and layout access methods & Acceptance check UC13 \\
FR15 & Statistics query & Acceptance check UC6 \\
FR16 & Stock consumed when an order completes & Unit test on the orderable list; acceptance check
FR16 \\
FR17 & Reservation derived from unfinished orders & Acceptance checks C2, both the holding and the
release \\
FR18 & Return queued on the same topic as deliveries & Exercised in the duty-state record; never
refused by construction \\
FR19 & Non-latching stop handled on the subscription & Acceptance checks UC11, including ground
truth that the robot did not move \\
FR20 & Session check on every page, role check on the request & Unit test on the gate across all
eight pages; acceptance checks UC12 \\
FR21 & Table name validated against the schema catalogue & Unit tests on the rejection and on
redaction \\
NFR1 & Timeout on the prediction client & Acceptance check NFR1, measured at 416 ms \\
NFR2 & Fallback path in the admission workflow & Unit test with an unreachable service \\
NFR3 & Server-side templates & Acceptance check NFR3 \\
NFR4 & One-second publication interval on the simulated clock & Largest gap between consecutive
telemetry rows over a delivery. \emph{Was not met when measured; see
Section~\ref{sec:nfr4}.} \\
NFR5 & Ground-truth comparison & Acceptance check UC5 \\
NFR6 & Generator and trainer kept outside the build & Static check that no module under the source
tree imports the training package \\
\bottomrule
\caption{Traceability from requirements to implementation and verification.}
\end{longtable}

\chapter{Figure Index}

\begin{longtable}{@{}p{0.10\textwidth}p{0.84\textwidth}@{}}
\toprule
\textbf{Figure} & \textbf{Content} \\
\midrule
\endhead
\ref{fig:warehouse} & Warehouse layout drawn to scale \\
\ref{fig:incremental} & The incremental model as followed, phases within each increment \\
\ref{fig:usecase} & Use case diagram \\
\ref{fig:domainclass} & Domain class diagram \\
\ref{fig:er} & Database schema, entity-relationship diagram \\
\ref{fig:dbclass} & Database schema, UML class diagram \\
\ref{fig:layers} & Layered architecture \\
\ref{fig:component} & Component diagram \\
\ref{fig:deployment} & Deployment diagram \\
\ref{fig:deps} & Dependency direction \\
\ref{fig:seqsuccess} & Scenario 1: an order placed and delivered \\
\ref{fig:seqrefused} & Scenario 2: an order refused on the robot's condition \\
\ref{fig:seqestop} & Scenario 3: a running order halted by the emergency stop \\
\ref{fig:payloadenergy} & Energy per metre against payload \\
\ref{fig:statemachine} & Duty-cycle state machine \\
\ref{fig:battery} & Battery over one session \\
\ref{fig:temperature} & Motor temperature over one session \\
\ref{fig:duty} & Duty states over one session \\
\ref{fig:wear} & Mechanical condition over one session \\
\ref{fig:activity} & Activity diagram of the admission workflow \\
\ref{fig:seqverify} & Sequence diagram of physical delivery verification (C3) \\
\ref{fig:mlpipeline} & Machine-learning pipeline \\
\ref{fig:mlbalance} & Class balance of the training set \\
\ref{fig:confusion} & Confusion matrix \\
\ref{fig:importance} & Feature importance \\
\ref{fig:pyramid} & Test strategy \\
\ref{fig:testcomp} & Composition of the test suite \\
\ref{fig:acceptance} & Acceptance results \\
\ref{fig:accuracy} & Delivery accuracy \\
\ref{fig:predvsactual} & Predicted against measured route length \\
\ref{fig:shotform} & Screenshot: the order form \\
\ref{fig:shotaccepted} & Screenshot: an accepted order preview \\
\ref{fig:shotrefused} & Screenshot: a refused order preview \\
\ref{fig:shotorders} & Screenshot: the order history \\
\ref{fig:shotrobot} & Screenshot: the robot condition page \\
\ref{fig:shotlogin} & Screenshot: the login page \\
\ref{fig:shotestop} & Screenshot: the order list after an emergency stop \\
\ref{fig:shotadminproducts} & Screenshot: product maintenance \\
\ref{fig:shotadmindb} & Screenshot: the read-only database browser \\
\ref{fig:shotreserved} & Screenshot: a preview taken against reserved work \\
\ref{fig:shotstock} & Screenshot: a delivered product marked out of stock \\
\ref{fig:shotgazebo} & Screenshot: the warehouse in the simulator \\
\ref{fig:shotrviz} & Screenshot: the warehouse as the robot perceives it \\
\ref{fig:shotlaunch} & Screenshot: the launch console when the robot is ready \\
\bottomrule
\end{longtable}

\end{document}
